\chapter{Entropic OT, Sinkhorn, and Barycentric Projection}
\label{ch:bf-dpf-entropic-ot-sinkhorn}

The previous chapter fixed the deterministic OT equal-weighting baseline.  This
chapter explains the main differentiable computational route used in the OT lane:
entropic regularization, finite Sinkhorn iteration, barycentric projection, and
the distinction between the mathematical teacher object and the numerical object
actually differentiated in code.

This chapter has one dominant goal: a reader should be able to understand what
object finite Sinkhorn computes, how barycentric output is produced from it, where
approximation enters, and what diagnostics must be checked before the result is
used in a downstream filtering scalar.

\section{What changes when entropy is added}
\label{sec:bf-eot-what-changes}

The previous chapter fixed the unregularized deterministic baseline by choosing a
feasible coupling with minimal declared transport cost,
\begin{equation}
\label{eq:bf-eot-deterministic-baseline-repeat}
  \Pi_0^\star
  \in
  \arg\min_{\Pi\in\mathcal U(w,u)}
  \langle C,\Pi\rangle.
\end{equation}
Entropic OT changes that teacher object before any numerical approximation enters.
For a regularization level \(\varepsilon>0\), it replaces the linear program by
an objective that also penalizes overly concentrated couplings through the entropy
term.  In plain language, the plan is now rewarded for spreading mass more smoothly
instead of concentrating all mass as sharply as possible.

This change matters mathematically as well as numerically.  The transport object is
no longer the unregularized deterministic optimizer from
\eqref{eq:bf-eot-deterministic-baseline-repeat}.  It is a different teacher problem
whose solution is smoother and whose first-order conditions can be rewritten as a
matrix-scaling problem.  That rewrite is what later makes Sinkhorn iteration
possible.  Thus the chapter's logic is:
\begin{enumerate}[label=(\roman*)]
  \item change the deterministic teacher by adding entropy,
  \item derive the factorized form of the new optimizer,
  \item derive the balancing equations for its scaling vectors,
  \item and only then discuss the finite numerical route used in code.
\end{enumerate}

The gain is numerical smoothness and a tractable scaling-based solver.  The cost is
that the transport object is no longer the unregularized deterministic OT solution.
This is the chapter's first approximation layer, and later finite iteration will add
another one on top of it.

\section{Running three-particle cell with entropic smoothing}
\label{sec:bf-eot-running-cell}

Keep the same running cell as before:
\[
  x_1=-2,
  \qquad x_2=0,
  \qquad x_3=3,
  \qquad
  w=
  \left(
    \frac{1}{2},
    \frac{1}{3},
    \frac{1}{6}
  \right),
  \qquad
  u=
  \left(
    \frac{1}{3},
    \frac{1}{3},
    \frac{1}{3}
  \right),
\]
with quadratic cost matrix
\begin{equation}
\label{eq:bf-eot-running-cost}
  C
  =
  \begin{pmatrix}
    0 & 4 & 25 \\
    4 & 0 & 9 \\
    25 & 9 & 0
  \end{pmatrix}.
\end{equation}
For \(\varepsilon=1\), the Gibbs kernel is
\begin{equation}
\label{eq:bf-eot-running-kernel}
  K_1=
  \begin{pmatrix}
    1 & e^{-4} & e^{-25} \\
    e^{-4} & 1 & e^{-9} \\
    e^{-25} & e^{-9} & 1
  \end{pmatrix}.
\end{equation}
The entries are heavily concentrated near the diagonal, but not exactly zero away
from it.  That is a simple visual sign of what entropy has changed: the plan is no
longer forced into the sharpest possible assignment.

\section{Entropic OT teacher object}
\label{sec:bf-eot-teacher-object}

Let \(\mathcal U(w,u)\) be the feasible coupling set from the previous chapter.
The entropically regularized OT problem is \citep{cuturi2013sinkhorn,peyre2019computational}
\begin{equation}
\label{eq:bf-eot-primal}
  \Pi_\varepsilon^\star
  =
  \arg\min_{\Pi\in\mathcal U(w,u)}
  \left\{
    \langle \Pi,C\rangle
    + \varepsilon\sum_{i,j}\Pi_{ij}(\log\Pi_{ij}-1)
  \right\}.
\end{equation}
This chapter uses the same entropy convention as the earlier OT material, so the
minimizer is the exact teacher for the chosen regularized object, not for the
unregularized linear program and not for categorical resampling.

\paragraph{Reader checkpoint.}
At this stage, the chapter has changed the mathematical object once: from the
unregularized deterministic OT coupling to the entropically smoothed OT coupling.
Nothing numerical has happened yet.  The next step still stays at the exact-object
level: first identify the factorized teacher form, and only after that turn to the
finite executable route.

\subsection{Exact teacher versus exact engineering route}
\label{sec:bf-eot-teacher-versus-engineering}

a useful extra distinction before moving on to
finite Sinkhorn iterations.  There are two different senses in which a route may
still be called ``exact'' for the entropic OT object.

First, there is the mathematical teacher object
\(\Pi_\varepsilon^\star\) from \eqref{eq:bf-eot-primal}.  That is the exact
minimizer of the entropically regularized problem.  Second, there is an
engineering route that preserves that same entropic object while changing how the
matrix-scaling updates are executed in memory or on hardware.  Exact online,
streaming, and GPU Sinkhorn methods belong to this second category: they do not
change the entropic teacher problem, but they avoid materializing the full cost,
kernel, or coupling whenever possible \citep{feydy2019geomloss,charlier2021keops,ye2026flashsinkhorn,xiao2026fastsinkhorn}.

This distinction matters for the BayesFilter transport lane.  A method that
recomputes the same entropic scaling updates through blocked reductions or
streamed kernel application is still on the semantics-preserving reference lane.
By contrast, a method that replaces the kernel, truncates the coupling family, or
changes the transport objective has moved to a different approximation class even
if it is motivated by scalability.  The reader should therefore keep separate:
\begin{enumerate}[label=(\roman*)]
  \item the exact entropic teacher object,
  \item the finite numerical object produced by a declared Sinkhorn execution
    path,
  \item the engineering strategy used to realize that numerical path.
\end{enumerate}
The next sections fix the second object explicitly and later staged chapters will
discuss scalability routes that either preserve or alter it.

\subsection{The current BayesFilter bottleneck}
\label{sec:bf-eot-current-bottleneck}

The need for scalable OT in BayesFilter is not only a generic literature concern.
The current TensorFlow route already distinguishes dense and streamed execution.
Dense mode materializes the effective transport matrix and applies it to the
particle cloud directly.  Streamed mode avoids storing that full matrix, but it
still loops over row and column blocks and therefore still computes the same
underlying pairwise transport interactions.

For equal-weight barycentric output, the streamed route still forms costs of the
schematic form
\begin{equation}
\label{eq:bf-eot-current-cost}
  C_{ij}
  =
  \frac{1}{2}\|\tilde x_i-\tilde x_j\|_2^2,
\end{equation}
then log-transport weights of the schematic form
\begin{equation}
\label{eq:bf-eot-current-log-transport}
  \log T_{ij}
  =
  \frac{f_i+g_j-C_{ij}}{\varepsilon}
  - \ell_j
  + \log N
  + \log w_j,
\end{equation}
and finally blockwise transport application by accumulation,
\begin{equation}
\label{eq:bf-eot-current-streamed-application}
  z_i = \sum_j T_{ij}x_j.
\end{equation}
The gain is real: memory pressure is reduced because the full transport matrix need
not be kept in RAM or device memory.  But the all-pairs arithmetic burden remains.
Streaming therefore solves part of the storage problem without yet solving the full
large-\(N\), large-\(d_x\) interaction problem.  This is exactly why the derivation
below matters: once the teacher object has been rewritten as a scaling problem, one
can see clearly which later routes merely change execution strategy and which ones
change the transport object itself.

\section{Why the optimizer factorizes}
\label{sec:bf-eot-factorization}

The next step is to derive the multiplicative structure of the entropic optimizer
rather than only state it.  Introduce Lagrange multipliers \(\alpha\in\mathbb R^N\)
for the row constraints and \(\beta\in\mathbb R^N\) for the column constraints.
With the entropy convention from \eqref{eq:bf-eot-primal}, the Lagrangian is
\begin{align}
\label{eq:bf-eot-lagrangian}
  \mathcal L(\Pi,\alpha,\beta)
  &=
  \sum_{i,j}
  \left[
    C_{ij}\Pi_{ij}
    + \varepsilon \Pi_{ij}(\log\Pi_{ij}-1)
  \right]
  + \sum_i \alpha_i\left(\sum_j\Pi_{ij}-w_i\right) \\
  &\qquad
  + \sum_j \beta_j\left(\sum_i\Pi_{ij}-u_j\right).
\end{align}
Stationarity with respect to each entry \(\Pi_{ij}\) gives
\begin{equation}
\label{eq:bf-eot-stationarity}
  0
  =
  \frac{\partial \mathcal L}{\partial \Pi_{ij}}
  =
  C_{ij}
  + \varepsilon \log \Pi_{ij}
  + \alpha_i + \beta_j,
\end{equation}
so
\begin{equation}
\label{eq:bf-eot-log-pi}
  \log\Pi_{ij}
  =
  -\frac{C_{ij}}{\varepsilon}
  -\frac{\alpha_i}{\varepsilon}
  -\frac{\beta_j}{\varepsilon}.
\end{equation}
Exponentiating entrywise yields
\begin{equation}
\label{eq:bf-eot-entrywise-factorization}
  \Pi_{ij}
  =
  \exp\!\left(-\frac{\alpha_i}{\varepsilon}\right)
  \exp\!\left(-\frac{C_{ij}}{\varepsilon}\right)
  \exp\!\left(-\frac{\beta_j}{\varepsilon}\right).
\end{equation}
Define the positive scalings
\begin{equation}
\label{eq:bf-eot-scaling-definitions}
  a_i = \exp(-\alpha_i/\varepsilon),
  \qquad
  b_j = \exp(-\beta_j/\varepsilon),
  \qquad
  (K_\varepsilon)_{ij}=\exp(-C_{ij}/\varepsilon).
\end{equation}
Then the optimizer must have the Gibbs-scaled form
\begin{equation}
\label{eq:bf-eot-factorization}
  \Pi_\varepsilon^\star
  =
  \diag(a)K_\varepsilon\diag(b).
\end{equation}
This is the key structural step.  The entropic teacher problem has been converted
from a constrained optimization over all couplings into the search for two positive
scaling vectors that make a known kernel satisfy the marginal constraints.

\subsection{Why the balancing equations are Sinkhorn updates}
\label{sec:bf-eot-balancing-equations}

Once \eqref{eq:bf-eot-factorization} is available, the marginal constraints from
\(\mathcal U(w,u)\) become explicit equations for the scaling vectors.  The row
constraints require
\begin{equation}
\label{eq:bf-eot-row-balance-derived}
  \Pi_\varepsilon^\star\mathbf 1
  =
  \diag(a)K_\varepsilon\diag(b)\mathbf 1
  =
  a \odot (K_\varepsilon b)
  = w,
\end{equation}
where \(\odot\) denotes entrywise multiplication.  Therefore
\begin{equation}
\label{eq:bf-eot-a-balance}
  a = w \oslash (K_\varepsilon b).
\end{equation}
Likewise the column constraints require
\begin{equation}
\label{eq:bf-eot-col-balance-derived}
  (\Pi_\varepsilon^\star)^\top\mathbf 1
  =
  \diag(b)K_\varepsilon^\top\diag(a)\mathbf 1
  =
  b \odot (K_\varepsilon^\top a)
  = u,
\end{equation}
so
\begin{equation}
\label{eq:bf-eot-b-balance}
  b = u \oslash (K_\varepsilon^\top a).
\end{equation}
These are exactly the balancing equations behind Sinkhorn iteration.  The algorithm
is not a mysterious OT oracle added from outside the derivation.  It simply
alternates the two marginal repairs implied by the factorized optimizer.

\paragraph{Reader checkpoint.}
The chapter has now introduced three different objects that should not be
confused:
\begin{enumerate}[label=(\roman*)]
  \item the exact entropic teacher problem \eqref{eq:bf-eot-primal},
  \item the factorized teacher optimizer \eqref{eq:bf-eot-factorization},
  \item and the balancing equations \eqref{eq:bf-eot-a-balance}--\eqref{eq:bf-eot-b-balance}
  that a finite numerical routine will try to satisfy.
\end{enumerate}
Finite Sinkhorn iteration begins only after these exact relations have been fixed.

\section{Algorithm: finite Sinkhorn equal-weighting}
\label{sec:bf-eot-algorithm}

\begin{algorithm}[ht]
\caption{\texttt{sinkhorn\_equalweight\_step}}
\begin{algorithmic}[1]
\Require weighted particle cloud \(\{x_i,w_i\}_{i=1}^N\), equal target marginal \(u\), cost matrix \(C\), regularization \(\varepsilon>0\), initialization \((a^{(0)},b^{(0)})\), iteration budget \(K\), and stabilization policy
\Ensure finite coupling \(\Pi_{\varepsilon,K}\), barycentric equal-weight cloud \(\widetilde X\), row/column residuals, and declared numerical metadata
\State Form the Gibbs kernel \((K_\varepsilon)_{ij}=\exp(-C_{ij}/\varepsilon)\)
\For{\(k=0,\ldots,K-1\)}
  \State Update \(a^{(k+1)} = w \oslash (K_\varepsilon b^{(k)})\)
  \State Update \(b^{(k+1)} = u \oslash (K_\varepsilon' a^{(k+1)})\)
\EndFor
\State Form the finite coupling \(\Pi_{\varepsilon,K}=\diag(a^{(K)})K_\varepsilon\diag(b^{(K)})\)
\State Compute residuals
\[
  r_{\mathrm{row}}=\|\Pi_{\varepsilon,K}\mathbf 1-w\|_1,
  \qquad
  r_{\mathrm{col}}=\|\Pi_{\varepsilon,K}'\mathbf 1-u\|_1
\]
\For{each output slot \(j=1,\ldots,N\)}
  \State Compute \(\tilde x^{(j)} = \frac{1}{u_j}\sum_{i=1}^N \Pi_{ij}x_i\)
\EndFor
\State Return the coupling, equal-weight cloud, residuals, \(\varepsilon\), \(K\), initialization policy, and stabilization metadata
\end{algorithmic}
\end{algorithm}

The algorithm returns a \emph{finite} numerical object, not the exact teacher
optimizer from \eqref{eq:bf-eot-primal}.  After \(K\) balancing passes, the
executed coupling is
\begin{equation}
\label{eq:bf-eot-finite-coupling-repeat}
  \Pi_{\varepsilon,K}
  =
  \diag(a^{(K)})K_\varepsilon\diag(b^{(K)}),
\end{equation}
with residuals that measure how closely the declared row and column constraints
have been met.  This distinction is the second approximation layer of the chapter:
first the deterministic teacher was replaced by an entropic teacher, and then that
entropic teacher was approximated by a finite numerical route.  The next custom-
gradient chapter will differentiate this declared finite route, not an abstract
teacher floating above the executed computation.

\paragraph{Reader checkpoint.}
It is useful to keep the chain of objects explicit:
\begin{enumerate}[label=(\roman*)]
  \item deterministic OT teacher \(\Pi_0^\star\),
  \item exact entropic teacher \(\Pi_\varepsilon^\star\),
  \item finite executed coupling \(\Pi_{\varepsilon,K}\),
  \item barycentric equal-weight cloud built from \(\Pi_{\varepsilon,K}\).
\end{enumerate}
Later differentiation claims must attach to the actually executed object in this
chain, not jump back silently to one of the earlier teachers.

\subsection{Exact online and streaming Sinkhorn as the reference lane}
\label{sec:bf-eot-reference-lane}

The scalable-OT survey also clarifies how BayesFilter should talk about scaling
without silently changing the transport object.  A finite Sinkhorn route can be
made substantially more practical by changing the execution strategy rather than
the underlying entropic teacher problem.  Streaming, online, or GPU-specialized
implementations still belong to the reference lane when they preserve the same
entropic scaling equations and the same barycentric output rule while avoiding
full materialization of the cost or kernel matrix \citep{feydy2019geomloss,charlier2021keops,ye2026flashsinkhorn,xiao2026fastsinkhorn}.

In that reference-lane setting, the implementation may compute the needed
row-wise or column-wise reductions in blocks, keep only partial transport state
in memory, or fuse barycentric application with solver passes.  What must remain
fixed is the declared numerical object: the finite coupling
\(\Pi_{\varepsilon,K}\), its residual checks, and the barycentric output derived
from that same finite coupling.  Thus a streamed or GPU-tiled Sinkhorn routine is
not automatically a new approximation class.  It becomes a new approximation
class only when the scaling equations, kernel, feasible family, or output object
itself are changed.

This exact-reference reading is important because it separates engineering success
from scientific overclaim.  Exact-online or GPU-specialized Sinkhorn may reduce
memory traffic, improve device utilization, and make larger particle clouds
practical while still computing the same finite entropic object.  But these routes
normally leave the dominant all-pairs interaction structure in place: every full
row or column update still depends on interactions with the whole opposing cloud.
So they are the right first scaling lane when the goal is to preserve semantics,
but they are not by themselves the final answer to truly large-\(N\), large-\(d_x\)
transport.

This chapter therefore treats exact-online and streaming implementations as
engineering variants of the same entropic reference lane.  Later staged chapters
on retained-teacher approximations will take up the different case where the
transport object itself is compressed, factorized, or otherwise altered for
scalability.

This is the chapter’s main callable procedure.  A reader who implements only one
OT route from the monograph should be able to implement this one.

\subsection{Sparse, screened, stabilized, and accelerated Sinkhorn variants}
\label{sec:bf-eot-solver-variant-boundary}

The scalable-OT survey also suggests a useful boundary inside the broader
Sinkhorn family.  Some variants still belong naturally to the exact-reference
lane because they mainly stabilize or schedule the same scaling equations more
carefully.  Others move toward approximation because they restrict which entries,
rows, or columns are treated as active, or because they replace dense kernel use
by a screened or sparse surrogate \citep{schmitzer2016sparse,alaya2019screenkhorn,schmitzer2019stabilized,altschuler2017nearlinear,dvurechensky2018computational,abid2018greedy,lin2019efficient,xu2026accelerating,genevay2016stochastic}.

At the most conservative end, stabilization policies remain part of the same
finite-Sinkhorn object.  Log-domain updates, shifting, epsilon scheduling, or
other numerical defenses are execution details attached to the same declared
coupling \(\Pi_{\varepsilon,K}\) and should therefore be judged first by the
usual marginal residual and barycentric parity checks.

Sparse or screened routes require a sharper audit.  If they only prune obviously
negligible work while preserving the same effective scaling equations and output
object to numerical tolerance, they may still be read as near-reference
engineering variants.  But once the active support, feasible family, or effective
kernel semantics are altered in a way that changes the finite transport object,
the method is no longer only a better implementation of the same reference lane.
It has crossed into retained-teacher or structure-changing territory and must be
judged accordingly.

A similar caution applies to accelerated, greedy, or stochastic update schedules.
These methods may reduce iteration count, repair the most violated row or column
first, or otherwise improve convergence behavior of the balancing routine.  That
is valuable, but it does not by itself remove dense-kernel access.  If each major
update still needs interactions against the whole opposing cloud, the dominant
large-\(N\), large-\(d_x\) arithmetic bottleneck remains.

This is why the present chapter keeps these variants conceptually downstream from
the finite teacher object rather than promoting them to new transport doctrines.
They refine how the reference route is executed.  They do not automatically change
what object the filter is claiming to compute.

\section{Worked balancing pass on the running cell}
\label{sec:bf-eot-worked-balance}

Take the simplest initialization \(b^{(0)}=(1,1,1)'\).  Then the first
row-balancing step computes
\begin{equation}
\label{eq:bf-eot-first-row-balance}
  a^{(1)} = w \oslash (K_1 b^{(0)}).
\end{equation}
For the running cell,
\[
  K_1 b^{(0)}
  \approx
  (1.0183,\,1.0184,\,1.0001),
\]
so
\[
  a^{(1)}
  \approx
  (0.491,\,0.327,\,0.167).
\]
The next column-balancing step then uses these row-repaired values to push the
column sums toward \((1/3,1/3,1/3)\).  The exact arithmetic is less important
than the interpretation: finite Sinkhorn repeatedly repairs whichever marginal is
currently wrong.

\paragraph{Plain-language takeaway.}
Sinkhorn is not a mysterious OT oracle.  It is an alternating balancing procedure
for a smoothed transport matrix.  That is why later learned warm-start methods try
to predict better initial balancing coordinates rather than replacing the solver
outright.

\section{Barycentric projection as the output map}
\label{sec:bf-eot-barycentric}

The solver output is still only a coupling matrix.  The equal-weight particle cloud
that later code consumes is obtained by barycentric projection:
\begin{equation}
\label{eq:bf-eot-barycentric}
  \tilde x^{(j)}
  =
  \frac{1}{u_j}\sum_{i=1}^N \Pi_{ij}x_i.
\end{equation}
In the equal-weight case, this is
\begin{equation}
\label{eq:bf-eot-barycentric-equal}
  \tilde x^{(j)}=N\sum_{i=1}^N \Pi_{ij}x_i.
\end{equation}
This step is where the numerical coupling becomes the actual equal-weight output
cloud.  It is therefore the place where an implementation must be especially clear
about shapes, normalization, and any auxiliary state that is or is not carried
alongside the particle locations.

\begin{proposition}[Barycentric equal-weighting contracts covariance]
\label{prop:bf-eot-barycentric-covariance-contraction}
Let \(X\) be a random variable that takes values \(x_1,\ldots,x_N\in\mathbb R^d\)
with probabilities \(w_1,\ldots,w_N\).  Let \(J\) be an equal-weight output index,
\(\mathbb P(J=j)=1/N\), and suppose a coupling between \(X\) and \(J\) has the
same source and target marginals as the OT coupling used in
\eqref{eq:bf-eot-barycentric-equal}.  The barycentric output associated with
index \(J\) is
\[
  \widetilde X = \mathbb E[X\mid J].
\]
Then
\begin{equation}
\label{eq:bf-eot-barycentric-mean-preservation}
  \mathbb E[\widetilde X]=\mathbb E[X],
\end{equation}
but
\begin{equation}
\label{eq:bf-eot-barycentric-covariance-contraction}
  \operatorname{Cov}(\widetilde X)
  =
  \operatorname{Cov}(X)
  -
  \mathbb E\!\left[\operatorname{Cov}(X\mid J)\right]
  \preceq
  \operatorname{Cov}(X).
\end{equation}
Equality holds only when the conditional covariance
\(\operatorname{Cov}(X\mid J)\) is zero almost surely.
\end{proposition}

\begin{proof}
The mean identity follows from the tower property:
\[
  \mathbb E[\widetilde X]
  =
  \mathbb E[\mathbb E[X\mid J]]
  =
  \mathbb E[X].
\]
For the covariance, apply the law of total covariance:
\[
  \operatorname{Cov}(X)
  =
  \operatorname{Cov}\!\left(\mathbb E[X\mid J]\right)
  +
  \mathbb E\!\left[\operatorname{Cov}(X\mid J)\right].
\]
Since \(\widetilde X=\mathbb E[X\mid J]\), rearranging gives
\eqref{eq:bf-eot-barycentric-covariance-contraction}.  The second term is
positive semidefinite, so the barycentric covariance is no larger in Loewner
order.  Equality requires the conditional covariance term to vanish almost
surely.
\end{proof}

\paragraph{Implication for filtering values.}
This proposition is the finite-\(N\) reason that a barycentric equal-weight cloud
should not be treated as a categorical-resampling cloud.  Categorical resampling
draws particles from the weighted empirical law and is unbiased for test
functions conditional on the weighted cloud.  Barycentric projection instead
replaces each conditional cloud by its conditional mean.  It preserves affine
test functions, including the first moment, but it generally changes quadratic
and nonlinear test functions.  In a linear-Gaussian state-space model, the next
predictive likelihood increment depends on the predictive covariance through
\(S_t=H P_{t\mid t-1}H'+R\).  A covariance-contracted reset can therefore change
the finite-\(N\) scalar being estimated even when Sinkhorn marginal residuals are
small, except in special cases where the lost covariance lies in observation-null
directions or otherwise cancels from the downstream likelihood increment.  The
exact Kalman likelihood remains a valid comparator for a no-OT weighted particle
approximation, but it is not a direct oracle for the barycentric-reset scalar
unless the reset route also preserves the relevant moments or the comparison is
explicitly redefined.

\paragraph{Implication for SV, SIR, and nonlinear models.}
The same issue is not special to LGSSM.  In stochastic-volatility models, future
observation densities are nonlinear functions of latent volatility, so the
finite-\(N\) likelihood scalar can change when barycentric projection compresses
the spread of the volatility particles.  In SIR and other nonlinear compartmental
models, the transition and observation maps are nonlinear in the state, so
preserving only the particle mean is not enough to preserve infection tails,
state constraints, threshold behavior, or future likelihood curvature.  A value
or gradient computed through the barycentric-reset filter can still be a correct
gradient of the executed barycentric scalar; it is not automatically the gradient
of the categorical-resampling particle-filter likelihood, the exact filtering
likelihood, or a moment-preserving approximation.  Consequently, SIR or SV
gradient checks must be same-scalar checks unless an additional evidence gate
shows that barycentric bias is negligible for the declared model, particle count,
regularization, and horizon.

The distinction matters for the later custom-gradient chapter
(Chapter~\ref{ch:bf-ledh-pfpf-ot-custom-gradient}).  That chapter does not
differentiate an abstract teacher floating outside the executed computation.  It
must differentiate the declared finite Sinkhorn route together with the barycentric
map that actually produces the equal-weight cloud used by the filtering scalar.

\section{Remedy options for barycentric covariance loss}
\label{sec:bf-eot-barycentric-remedies}

The remedy has to be stated at the level of the reset map, because the
covariance defect is not a vague numerical pathology of Sinkhorn.  It is a
deterministic consequence of replacing each column of a transport plan by its
conditional mean.  This subsection makes that statement as a finite-dimensional
matrix calculation.

For the algebra in this subsection, assume first that the Sinkhorn marginal
constraints have been solved exactly.  Finite iteration and floating-point
residuals add normalization errors that should be checked separately, but they
are not the mechanism discussed here.  Let
\[
  X = [x_1,\ldots,x_N]\in\mathbb R^{d_x\times N},
  \qquad
  w\in\mathbb R^N_+,\qquad
  \mathbf 1^\top w=1,
  \qquad
  W=\diag(w).
\]
The weighted empirical measure before reset is
\[
  \nu_w=\sum_{i=1}^N w_i\delta_{x_i}.
\]
Any deterministic equal-weight reset whose output particles are linear
combinations of the input particles can be written as
\begin{equation}
\label{eq:bf-eot-reset-transform}
  Y = XD,
  \qquad
  y_j=\sum_{i=1}^N x_i D_{ij},
\end{equation}
for some \(N\times N\) transform matrix \(D\).  The barycentric Sinkhorn reset is
the special case
\begin{equation}
\label{eq:bf-eot-barycentric-D}
  D_B = N\Pi_{\varepsilon,K},
\end{equation}
because the target marginal is \(u_j=1/N\).  The row and column marginal
constraints on \(\Pi_{\varepsilon,K}\) become the first-order transform
constraints
\begin{equation}
\label{eq:bf-eot-first-order-transform-contract}
  D^\top \mathbf 1 = \mathbf 1,
  \qquad
  D\mathbf 1 = Nw.
\end{equation}
The first condition says that every output particle is an affine combination of
the input particles.  The second says that the equal-weight output cloud has the
same mean as the weighted input cloud:
\begin{equation}
\label{eq:bf-eot-transform-mean}
  \bar y_D
  :=
  \frac1N Y\mathbf 1
  =
  \frac1N XD\mathbf 1
  =
  Xw
  =:
  \mu_w.
\end{equation}

The second moment is where the problem lives.  Define the weighted input
covariance and the equal-weight output covariance by
\begin{align}
\label{eq:bf-eot-weighted-covariance}
  \Sigma_w
  &:=
  \sum_{i=1}^N w_i(x_i-\mu_w)(x_i-\mu_w)^\top
  =
  X(W-ww^\top)X^\top,\\
\label{eq:bf-eot-transform-covariance}
  \Sigma_D
  &:=
  \frac1N\sum_{j=1}^N (y_j-\mu_w)(y_j-\mu_w)^\top
  =
  \frac1N X(D-w\mathbf 1^\top)(D-w\mathbf 1^\top)^\top X^\top.
\end{align}
A deterministic reset is therefore \emph{second-order exact for the current
cloud} when it satisfies \eqref{eq:bf-eot-first-order-transform-contract} and
\begin{equation}
\label{eq:bf-eot-second-order-cloud-contract}
  \Sigma_D=\Sigma_w.
\end{equation}
A stronger, cloud-independent sufficient condition is
\begin{equation}
\label{eq:bf-eot-second-order-matrix-contract}
  (D-w\mathbf 1^\top)(D-w\mathbf 1^\top)^\top
  =
  N(W-ww^\top).
\end{equation}
Equations
\eqref{eq:bf-eot-first-order-transform-contract}--\eqref{eq:bf-eot-second-order-matrix-contract}
are the finite-dimensional version of the second-order accurate ensemble-transform
condition used by Acevedo, de Wiljes, and Reich: the output ensemble matches both
the importance-weighted mean and covariance; see their Definition~3.1 and
Sections~5--6 for the Sinkhorn-based second-order ETPF construction
\citep{acevedodewiljesreich2017secondorder}.
The factor \(N\) in \eqref{eq:bf-eot-second-order-matrix-contract} is not a
convention-free decoration: it follows from using the biased equal-weight output
covariance \(N^{-1}\sum_j(y_j-\mu_w)(y_j-\mu_w)^\top\) together with a transform
matrix satisfying \(D\mathbf 1=Nw\).

\begin{proposition}[Why the barycentric Sinkhorn reset is not second-order exact]
\label{prop:bf-eot-barycentric-not-second-order}
Assume \(D_B=N\Pi_{\varepsilon,K}\) is nonnegative and satisfies
\eqref{eq:bf-eot-first-order-transform-contract}.  Interpret \(J\) as a uniform
output index and define a source index \(I\) by
\[
  \mathbb P(I=i\mid J=j)=D_{B,ij}.
\]
Then \(X_I\) has the weighted empirical law \(w\), the barycentric output is
\(Y_J=\mathbb E[X_I\mid J]\), and
\begin{equation}
\label{eq:bf-eot-D-covariance-gap}
  \Sigma_w-\Sigma_{D_B}
  =
  \mathbb E\!\left[\operatorname{Cov}(X_I\mid J)\right]
  \succeq 0.
\end{equation}
Consequently \(D_B\) satisfies the second-order cloud contract
\eqref{eq:bf-eot-second-order-cloud-contract} only if the conditional covariance
inside every output column vanishes as a state-space covariance.  Equivalently,
for each output column \(j\), all source particles with \(D_{B,ij}>0\) must have
the same state value, up to duplicate particle locations.  In generic diffuse
entropic couplings this condition fails.
\end{proposition}

\begin{proof}
The column sums in \eqref{eq:bf-eot-first-order-transform-contract} make
\(D_{B,\cdot j}\) a probability vector.  Since
\((1/N)D_B\mathbf 1=w\), the unconditional law of \(I\) is \(w\).  The output
particle in column \(j\) is exactly
\[
  y_j=\sum_i x_iD_{B,ij}=\mathbb E[X_I\mid J=j].
\]
Applying the law of total covariance gives
\[
  \operatorname{Cov}(X_I)
  =
  \operatorname{Cov}(\mathbb E[X_I\mid J])
  +
  \mathbb E[\operatorname{Cov}(X_I\mid J)].
\]
Because \(J\) is uniform,
\[
  \operatorname{Cov}(\mathbb E[X_I\mid J])
  =
  \frac1N\sum_{j=1}^N (y_j-\mu_w)(y_j-\mu_w)^\top
  =
  \Sigma_{D_B}.
\]
The left side is \(\Sigma_w\), which proves \eqref{eq:bf-eot-D-covariance-gap}.
\end{proof}

This is why the bug is not solved by asking Sinkhorn to converge more tightly.
A smaller marginal residual makes \(D_B\) a better approximation to the declared
entropic barycentric transform.  It does not change the identity
\eqref{eq:bf-eot-D-covariance-gap}, because that identity is about the output
map \(D_B\mapsto XD_B\), not about whether the coupling has been balanced
accurately.  A real remedy must therefore change, or explicitly keep, the
mathematical object being estimated.  The following contracts are mutually
different finite-\(N\) estimators.

\paragraph{Contract A: keep the same barycentric scalar.}
Set \(D=D_B\) and accept that \(\Sigma_D\preceq\Sigma_w\).  The reported value and
gradient are then only claims about the finite barycentric-reset scalar.  The
correct derivative check is a same-scalar finite-difference check of that exact
executed computation.  This contract repairs value-gradient naming, not the
covariance loss.  In notation, if the filter returns \(L_B(\theta)\), where every
reset uses \(Y_B(\theta)=X(\theta)D_B(\theta)\), then the admissible gradient
claim is
\[
  \nabla_\theta L_B(\theta),
\]
not the gradient of the categorical particle-filter likelihood and not the
gradient of an exact Kalman or nonlinear filtering likelihood.

\paragraph{Contract B: stochastic categorical output.}
Given the same nonnegative \(D_B\), draw \(I_j\) independently or by a chosen
resampling scheme from
\(\mathbb P(I_j=i\mid J=j)=D_{B,ij}\) and output \(z_j=x_{I_j}\).  Then for any
test function \(\varphi\) on the particle support,
\begin{equation}
\label{eq:bf-eot-stochastic-unbiased-test}
  \mathbb E\!\left[
    \frac1N\sum_{j=1}^N\varphi(z_j)\,\middle|\,D_B,X
  \right]
  =
  \sum_{i=1}^N w_i\varphi(x_i).
\end{equation}
This restores the unbiased empirical-measure property of categorical resampling
in expectation
\citep{gordon1993novel,doucet2001sequential,chopin2020introduction}, but the
realized reset is random and discontinuous as a function of weights and particles.
It is therefore a different estimator and not a pathwise differentiable
replacement for \(D_B\).

\paragraph{Contract C: second-order deterministic transform.}
Replace \(D_B\) by
\begin{equation}
\label{eq:bf-eot-second-order-correction}
  D^\sharp = D_B+\Delta,
  \qquad
  \Delta^\top\mathbf 1=0,
  \qquad
  \Delta\mathbf 1=0,
\end{equation}
and choose \(\Delta\) so that
\begin{equation}
\label{eq:bf-eot-second-order-correction-cov}
  \frac1N X(D^\sharp-w\mathbf 1^\top)
  (D^\sharp-w\mathbf 1^\top)^\top X^\top
  =
  X(W-ww^\top)X^\top.
\end{equation}
The stronger matrix-level target is
\begin{equation}
\label{eq:bf-eot-second-order-correction-matrix}
  (D^\sharp-w\mathbf 1^\top)(D^\sharp-w\mathbf 1^\top)^\top
  =
  N(W-ww^\top).
\end{equation}
One concrete reference object satisfying the matrix target is the NETF-type
square-root transform
\begin{equation}
\label{eq:bf-eot-netf-reference-transform}
  D_{\mathrm{NETF}}(Q)
  =
  w\mathbf 1^\top
  +
  \sqrt{N}\,(W-ww^\top)^{1/2}Q,
  \qquad
  Q\in\mathbb R^{N\times N},
  \qquad
  Q^\top Q=QQ^\top=I_N,
  \qquad
  Q\mathbf 1=\mathbf 1.
\end{equation}
Here \((W-ww^\top)^{1/2}\) denotes the principal symmetric square root.  Indeed,
\((W-ww^\top)\mathbf 1=0\), so the square root also annihilates \(\mathbf 1\), and
\(D_{\mathrm{NETF}}(Q)^\top\mathbf 1=\mathbf 1\),
\(D_{\mathrm{NETF}}(Q)\mathbf 1=Nw\), and
\eqref{eq:bf-eot-second-order-matrix-contract} holds:
\[
  (D_{\mathrm{NETF}}-w\mathbf 1^\top)
  (D_{\mathrm{NETF}}-w\mathbf 1^\top)^\top
  =
  N(W-ww^\top)^{1/2}QQ^\top(W-ww^\top)^{1/2}
  =
  N(W-ww^\top).
\]
Thus Contract C is mathematically explicit: the filter is no longer using the
barycentric transform \(D_B\); it is using a new transform \(D^\sharp\) whose
mean and covariance are the declared finite-\(N\) moments.  Inside this contract,
the NETF-type transform above and a second-order corrected ETPF are distinct
finite-\(N\) choices.  The NETF-type transform is a reference moment-preserving
object, not a promise of closeness to \(D_B\).  A second-order corrected ETPF
instead tries to stay close to the OT transform while enforcing the same moment
equations
\citep{reich2013nonparametric,acevedodewiljesreich2017secondorder}.
The unavoidable cost is that \(D^\sharp\) need not be nonnegative.  Hence the
corrected particles need not lie in the convex hull of the pre-reset particles,
and bounded, positive, simplex-valued, or otherwise constrained state spaces need
an explicit constraint policy.

\paragraph{Contract D: affine state-space moment restoration.}
Instead of correcting the \(N\times N\) transform matrix, one may correct the
already formed barycentric cloud.  Let \(Y_B=XD_B\) and
\(\Sigma_B=\Sigma_{D_B}\).  If there exists a matrix \(A_m\in\mathbb R^{d_x\times d_x}\)
such that
\begin{equation}
\label{eq:bf-eot-affine-moment-map}
  A_m\Sigma_B A_m^\top = \Sigma_w,
\end{equation}
then
\begin{equation}
\label{eq:bf-eot-affine-moment-output}
  y_j^\sharp
  =
  \mu_w + A_m(y_j-\mu_w),
  \qquad j=1,\ldots,N,
\end{equation}
has equal-weight mean \(\mu_w\) and covariance \(\Sigma_w\).  A full-rank example
is \(A_m=\Sigma_w^{1/2}\Sigma_B^{-1/2}\).  If \(\Sigma_B\) is singular, exact
affine restoration can exist only if
\(\operatorname{rank}(\Sigma_w)\le \operatorname{rank}(\Sigma_B)\);
jitter, truncation, or pseudoinverses define an additional approximation policy,
not an identity.  This contract is simple to test in an LGSSM, but it is a new
finite target and requires its own value gate, same-scalar gradient gate, and
state-constraint audit.

\paragraph{Contract E: positive spread, residual noise, and affine restoration.}
The preceding contracts leave a gap between two desirable but incompatible
properties.  Contract C gives an exact second-order deterministic transform, but
the transform can be signed.  Contract D gives a simple state-space covariance
repair, but it can be ill-conditioned when the barycentric cloud has already
collapsed in one or more directions.  A natural intermediate construction is to
split the job into three explicit stages:
\[
  D_B \longrightarrow D^+
  \longrightarrow \widetilde Y
  \longrightarrow Y^\star .
\]
Here \(D^+\) is still a nonnegative first-order transform, \(\widetilde Y\)
adds a small reparameterized residual in the missing covariance directions, and
\(Y^\star\) applies the final affine moment restoration.  This is not a
drop-in form of the second-order ETPF of
\citet{acevedodewiljesreich2017secondorder}, nor is it the nonlinear
transport-map framework of \citet{spantini2022coupling}.  It is a new finite
reset contract assembled from familiar components: positive ensemble transforms
\citep{reich2013nonparametric,corenflos2021differentiable}, and nearby classes
of second-order, stochastic, conditional, and learned ensemble transforms.  The
exact three-stage composition below was not found in the inspected sources, so
the surrounding design-space remarks are intentionally narrower than theorem- or
algorithm-level support.  The citations in this paragraph are therefore used only
where their source identity is already stable in the rewrite tree; unresolved or
provisional neighboring records are not needed for the mathematical contract
stated below and are omitted here.  The citations support ingredients and nearby
design space; they do not prove that this contract is stable, unbiased, or
preferable for LEDH-PFPF-OT.  Those are separate value, gradient,
conditioning, and constraint gates.

The construction is easiest to state on the linear support of the weighted
cloud.  Write
\[
  \mathcal S_w:=\operatorname{range}(\Sigma_w),
  \qquad
  P_w:=\hbox{the orthogonal projector onto }\mathcal S_w .
\]
Let \(r_w:=\operatorname{rank}(\Sigma_w)\).  Since \(\Sigma_w\) is the covariance
of \(N\) centered weighted particles, \(r_w\le N-1\); if a modified construction
asks for a rank larger than \(N-1\), no equal-weight cloud of \(N\) particles can
realize it.  All square roots below are principal square roots on their stated
supports.  If
the rank of \(\Sigma_w\), \(G_+\), or \(\widetilde\Sigma\) changes along a
parameter path, the corresponding square-root, pseudoinverse, or support
projector derivative may fail to be classical.  A differentiable implementation
therefore needs either a fixed-rank contract or an explicit spectral floor and
smoothed spectral projector policy.  This is a numerical assumption, not a
consequence of optimal transport.

The first stage chooses
\begin{equation}
\label{eq:bf-eot-positive-spread-transform}
  D^+\in\mathcal D_1^+
  :=
  \{D\in\mathbb R_+^{N\times N}:D^\top\mathbf 1=\mathbf 1,\,
  D\mathbf 1=Nw\}.
\end{equation}
The notation \(D^+\) does not mean that the matrix is invertible.  The relevant
rank condition is a state-space condition on the cloud
\[
  Y^+ := XD^+,
  \qquad
  \Sigma_+
  :=
  \frac1N
  (Y^+-\mu_w\mathbf 1^\top)(Y^+-\mu_w\mathbf 1^\top)^\top .
\]
The selection rule for \(D^+\) may be a sharper Sinkhorn plan, a spread-penalized
positive transport solve, or a positive projection of a partial second-order
correction.  The mathematical contract is weaker and more auditable than any
particular optimizer: \(D^+\) must remain in \(\mathcal D_1^+\), must be close
enough to the declared transport target for the application, and must pass a
spectral gate on \(\Sigma_+\).  If the gate fails, this contract has not produced
a usable reset.

\begin{proposition}[The residual covariance to be repaired is positive]
\label{prop:bf-eot-positive-spread-gap}
Let \(D^+\in\mathcal D_1^+\).  Define
\[
  G_+ := \Sigma_w-\Sigma_+ .
\]
Then
\begin{equation}
\label{eq:bf-eot-positive-spread-gap-formula}
  G_+
  =
  \frac1N
  \sum_{j=1}^N
  \sum_{i=1}^N
  D^+_{ij}
  (x_i-y^+_j)(x_i-y^+_j)^\top
  \succeq 0.
\end{equation}
Equality holds if and only if each output column of \(D^+\) has zero conditional
state covariance, in the same sense as
Proposition~\ref{prop:bf-eot-barycentric-not-second-order}.  Moreover
\(\operatorname{range}(G_+)\subseteq\operatorname{range}(\Sigma_w)\).
\end{proposition}

\begin{proof}
The proof is the same total-covariance argument as
Proposition~\ref{prop:bf-eot-barycentric-not-second-order}, with \(D_B\)
replaced by \(D^+\).  Nonnegativity and the column-sum constraint make each
column \(D^+_{\cdot j}\) a conditional probability vector.  The row-sum
constraint gives the unconditional law \(w\).  Therefore the covariance of the
weighted source draw decomposes into the covariance of the column barycentres
plus the mean conditional covariance inside the columns.  The first term is
\(\Sigma_+\), and the second is exactly the right-hand side of
\eqref{eq:bf-eot-positive-spread-gap-formula}.  This proves positive
semidefiniteness.  Since \(0\preceq G_+\preceq \Sigma_w\), every vector in
\(\ker(\Sigma_w)\) is also in \(\ker(G_+)\); for symmetric positive
semidefinite matrices this is equivalent to
\(\operatorname{range}(G_+)\subseteq\operatorname{range}(\Sigma_w)\).
\end{proof}

The second stage adds a differentiable residual in the directions represented by
\(G_+\).  Let \(\Xi\in\mathbb R^{d_x\times N}\) be an auxiliary base-noise
matrix, independent of the model parameters, satisfying
\[
  \mathbb E[\Xi]=0,
  \qquad
  \Xi\mathbf 1=0,
  \qquad
  \mathbb E\!\left[\frac1N\Xi\Xi^\top\right]=I_{d_x}.
\]
One concrete reparameterized choice is
\[
  H_N:=I_N-\frac1N\mathbf 1\mathbf 1^\top,
  \qquad
  \Xi:=\sqrt{\frac{N}{N-1}}\, ZH_N,
  \qquad
  Z_{ij}\stackrel{\mathrm{iid}}{\sim}N(0,1),
  \qquad N\ge2.
\]
The matrix \(Z\) is drawn independently of \(X\) and \(D^+\).  Then
\(\Xi\mathbf 1=0\), and
\(\mathbb E[N^{-1}\Xi\Xi^\top]=I_{d_x}\) because \(H_N\) is an idempotent
projector with trace \(N-1\).  A fixed quasi-random or orthogonal design may be
used instead, but then its empirical covariance is part of the finite contract
and must be reported.  For \(0\le\rho\le1\) and \(\tau\ge0\), define
\begin{equation}
\label{eq:bf-eot-residual-noise-cloud}
  \widetilde Y
  :=
  Y^+
  +
  B_{\rho,\tau}\Xi,
  \qquad
  B_{\rho,\tau}
  :=
  \sqrt{\rho}\,(G_+ + \tau P_w)^{1/2}.
\end{equation}
The parameter \(\rho\) controls how much of the gap is injected before the
affine map.  The parameter \(\tau\) is not a proof device; it is a numerical
floor on \(\mathcal S_w\).  Since
\(\operatorname{range}(G_+)\subseteq\mathcal S_w\), the residual covariance
\(G_+ + \tau P_w\) is positive semidefinite and has no component orthogonal to
the weighted cloud support.  If \(\tau\) is set too small, the affine whitening
in the next stage can amplify the auxiliary noise violently.  Therefore
\(\tau\) and \(\rho\) must be chosen by conditioning gates, not by the desire to
make the perturbation invisible.

\begin{proposition}[Expected effect of the residual stage]
\label{prop:bf-eot-residual-expected-covariance}
Assume \(D^+\in\mathcal D_1^+\), \(\mathbb E[\Xi\mid X,D^+]=0\),
\(\Xi\mathbf 1=0\), and
\(\mathbb E[N^{-1}\Xi\Xi^\top\mid X,D^+]=I_{d_x}\).  Conditional on \(X\) and \(D^+\),
the residual cloud \(\widetilde Y\) in
\eqref{eq:bf-eot-residual-noise-cloud} has mean \(\mu_w\), and its covariance
around \(\mu_w\) satisfies
\begin{equation}
\label{eq:bf-eot-residual-expected-covariance}
  \mathbb E\!\left[
    \frac1N(\widetilde Y-\mu_w\mathbf 1^\top)
    (\widetilde Y-\mu_w\mathbf 1^\top)^\top
    \,\middle|\,X,D^+
  \right]
  =
  \Sigma_+ + \rho(G_+ + \tau P_w).
\end{equation}
In particular, when \(\rho=1\) and \(\tau=0\), the conditional expected
covariance is \(\Sigma_w\).
\end{proposition}

\begin{proof}
Because \(D^+\mathbf 1=Nw\), the equal-weight mean of \(Y^+\) is \(\mu_w\).
Because \(\Xi\mathbf 1=0\), the equal-weight mean of \(B_{\rho,\tau}\Xi\) is
zero.  Expanding the covariance of \(\widetilde Y\) around \(\mu_w\) gives the
covariance of \(Y^+\), two cross terms, and the residual covariance.  The cross
terms have conditional expectation zero because the auxiliary noise has
conditional mean zero given \(X,D^+\).  The residual covariance has conditional
expectation
\[
  B_{\rho,\tau}
  \mathbb E\!\left[\frac1N\Xi\Xi^\top\,\middle|\,X,D^+\right]
  B_{\rho,\tau}^\top
  =
  B_{\rho,\tau}B_{\rho,\tau}^\top
  =
  \rho(G_+ + \tau P_w),
\]
by the definition of the principal square root.  This proves
\eqref{eq:bf-eot-residual-expected-covariance}.
\end{proof}

The proposition is deliberately an expectation statement.  For a realized finite
noise matrix \(\Xi\), the covariance of \(\widetilde Y\) need not equal
\(\Sigma_w\) even when \(\rho=1\) and \(\tau=0\).  This is not a defect in the
algebra; it is the price of using a reparameterized stochastic residual rather
than a discontinuous categorical resampling step.  The next affine stage turns
the realized cloud into an exact finite second-moment object only when the
support condition in Proposition~\ref{prop:bf-eot-residual-affine-restoration}
holds and the exact Moore--Penrose inverse square root is used.

\begin{proposition}[A small residual can repair support, but not sample-rank limits]
\label{prop:bf-eot-residual-support-repair}
Assume \(r_w\le N-1\), \(0<\rho\le1\), and \(\tau>0\).  Draw \(Z\) with
entries having a joint density, for example independent standard Gaussian
entries, and set
\(\Xi=\sqrt{N/(N-1)}\,ZH_N\).  Conditional on \(X\) and \(D^+\), the residual
matrix \(B_{\rho,\tau}\Xi\) has column span equal to \(\mathcal S_w\) with
probability one.  Consequently
\[
  \operatorname{range}(\Sigma_w)
  \subseteq
  \operatorname{range}(\widetilde\Sigma)
\]
with probability one.  If \(r_w>N-1\), no choice of \(N\) equal-weight output
particles can realize covariance \(\Sigma_w\), because every centered
\(d_x\times N\) cloud has rank at most \(N-1\).
\end{proposition}

\begin{proof}
Since \(\tau>0\),
\[
  G_+ + \tau P_w
\]
is strictly positive definite on \(\mathcal S_w\) and zero on
\(\mathcal S_w^\perp\).  Hence \(B_{\rho,\tau}\) has range \(\mathcal S_w\) and
rank \(r_w\).  Let \(U\in\mathbb R^{d_x\times r_w}\) have orthonormal columns
spanning \(\mathcal S_w\).  The centering matrix \(H_N\) has rank \(N-1\), and
\(U^\top ZH_N\) has rank \(r_w\) almost surely when \(r_w\le N-1\), because its
\(r_w\times (N-1)\) coordinate representation on
\(\operatorname{range}(H_N)\) has a joint density and at least one
\(r_w\times r_w\) minor is a nonzero polynomial in those variables.
Multiplication by the nonsingular
restriction of \(B_{\rho,\tau}\) on \(\mathcal S_w\) therefore produces a
residual cloud whose span is exactly \(\mathcal S_w\) almost surely.

The deterministic centered cloud \(Y^+-\mu_w\mathbf 1^\top\) also lies in
\(\mathcal S_w\).  The event that adding this fixed matrix to the full-rank
random residual reduces the rank on \(\mathcal S_w\) is again a polynomial zero
event.  To see that the relevant polynomial is not identically zero, choose any
centered matrix \(R_0\) with column span \(\mathcal S_w\); then the same
\(r_w\times r_w\) minor of
\((Y^+-\mu_w\mathbf 1^\top)+tR_0\) is nonzero for all sufficiently large
\(|t|\).  Hence the zero set has probability zero under the joint-density
assumption.  Thus the realized centered cloud
\(\widetilde Y-\mu_w\mathbf 1^\top\) has rank \(r_w\) on \(\mathcal S_w\) almost
surely, which is the displayed support inclusion.  The final statement follows
because any centered cloud \(Y-\bar y\mathbf 1^\top\) has columns summing to zero
and therefore rank at most \(N-1\).
\end{proof}

The third stage removes the remaining finite-sample mismatch by affine
restoration.  Define the realized residual covariance
\[
  \widetilde\Sigma
  :=
  \frac1N
  (\widetilde Y-\mu_w\mathbf 1^\top)
  (\widetilde Y-\mu_w\mathbf 1^\top)^\top .
\]
If \(\widetilde\Sigma\) is nonsingular on the support of \(\Sigma_w\), set
\begin{equation}
\label{eq:bf-eot-residual-affine-map}
  A_{\rho,\tau}
  :=
  \Sigma_w^{1/2}\,\widetilde\Sigma^{\dagger,-1/2},
  \qquad
  Y^\star
  :=
  \mu_w\mathbf 1^\top
  +
  A_{\rho,\tau}
  (\widetilde Y-\mu_w\mathbf 1^\top),
\end{equation}
where \(\widetilde\Sigma^{\dagger,-1/2}\) denotes the Moore--Penrose inverse of
the principal square root.

\begin{proposition}[Affine restoration after the residual stage]
\label{prop:bf-eot-residual-affine-restoration}
Assume \(D^+\in\mathcal D_1^+\), \(\Xi\mathbf 1=0\), and
\[
  \operatorname{range}(\Sigma_w)
  \subseteq
  \operatorname{range}(\widetilde\Sigma).
\]
Then the output \(Y^\star\) in \eqref{eq:bf-eot-residual-affine-map} has
equal-weight mean \(\mu_w\) and equal-weight covariance \(\Sigma_w\).
\end{proposition}

\begin{proof}
The mean statement follows from \(D^+\mathbf 1=Nw\) and \(\Xi\mathbf 1=0\).
For the covariance, let \(P_{\widetilde\Sigma}\) be the orthogonal projector
onto \(\operatorname{range}(\widetilde\Sigma)\).  By the definition of the
Moore--Penrose inverse square root,
\[
  \widetilde\Sigma^{\dagger,-1/2}
  \widetilde\Sigma
  \widetilde\Sigma^{\dagger,-1/2}
  =
  P_{\widetilde\Sigma}.
\]
The range assumption implies that this projector acts as the identity on
\(\operatorname{range}(\Sigma_w)\).  Since the range of
\(\Sigma_w^{1/2}\) is \(\operatorname{range}(\Sigma_w)\), this also implies
\(P_{\widetilde\Sigma}\Sigma_w^{1/2}=\Sigma_w^{1/2}\).  By symmetry of
\(P_{\widetilde\Sigma}\) and \(\Sigma_w^{1/2}\), it also implies
\(\Sigma_w^{1/2}P_{\widetilde\Sigma}=\Sigma_w^{1/2}\).  Hence
\[
  \frac1N
  (Y^\star-\mu_w\mathbf 1^\top)
  (Y^\star-\mu_w\mathbf 1^\top)^\top
  =
  \Sigma_w^{1/2}
  P_{\widetilde\Sigma}
  \Sigma_w^{1/2}
  =
  \Sigma_w .
\]
\end{proof}

An implementation may replace \(\widetilde\Sigma^{\dagger,-1/2}\) by a floored
or otherwise regularized inverse square root for conditioning.  That defines an
approximate restoration policy, not Proposition~\ref{prop:bf-eot-residual-affine-restoration}.
The approximation must report its spectral floor, covariance residual, and
conditioning diagnostics separately from the exact algebra above.

Contract E is attractive because it does not ask the nonnegative transform to do
the impossible.  Positivity is used only to build a noncollapsed first-order
cloud; residual noise supplies controlled missing directions; the affine map
enforces the declared second moment.  It is also dangerous for exactly the same
reason.  The final particles need not remain in the convex hull of the source
particles after the residual and affine stages, and a small residual can be made
large by \(\widetilde\Sigma^{-1/2}\).  The contract is therefore admissible only
with explicit gates:
\[
  \operatorname{rank}_{\tau_{\mathrm{eig}}}(\widetilde\Sigma)=r_w
  \quad\hbox{or}\quad
  \|(I-P_{\widetilde\Sigma,\tau_{\mathrm{eig}}})P_w\|\le \epsilon_{\mathrm{supp}},
  \qquad
  \lambda_{\min}^{\tau}(\widetilde\Sigma)>\tau_{\mathrm{eig}},
  \qquad
  \kappa^{\tau}(\widetilde\Sigma)<\kappa_{\max},
  \qquad
  \frac{\|\Sigma_w-\Sigma_+\|}{\|\Sigma_w\|}\le c_{\mathrm{gap}}
  \quad\hbox{or an explained exception.}
\]
where \(P_{\widetilde\Sigma,\tau_{\mathrm{eig}}}\) denotes the numerical support
projector obtained from the same eigenvalue threshold used by the rank gate.
The spectral quantities
\(\lambda_{\min}^{\tau}(\widetilde\Sigma)\) and
\(\kappa^{\tau}(\widetilde\Sigma)\) are computed only on this retained
\(\tau_{\mathrm{eig}}\)-support.  Eigenvalues below the threshold are treated as
discarded numerical null directions and must be charged to the support-projector
gate above; they must not simultaneously veto the fallback through a full
positive-spectrum condition.  If the implementation claims the exact
Moore--Penrose restoration rather than a floored numerical restoration, then the
full positive spectrum must be used and the support-projector fallback is no
longer sufficient.
It must also record the auxiliary-noise law or fixed design, random seeds,
Monte Carlo standard errors, same-scalar finite-difference checks, and
constraint violations after the affine map.  In SIR, SV, and other constrained
models, the construction should normally be tested first in validated
unconstrained or transformed coordinates.

\paragraph{Contract E--Chol: ridged square-root gauge for differentiability.}
The support-exact form above is mathematically natural but a poor default
coordinate system for a manual gradient implementation.  It asks the code to
differentiate eigenvectors, support projectors, principal square roots, and
Moore--Penrose inverse square roots.  When eigenvalues repeat, eigenvectors are
not unique.  When an eigenvalue crosses the numerical floor, the retained rank
changes.  Thus an eigendecomposition implementation can be perfectly reasonable
for a forward diagnostic while still being a bad coordinate chart for a
gradient-bearing reset.  This is the same engineering lesson as the
square-root sigma-point discussion in Chapter~\ref{ch:bf-svd-sigma-point}: the
covariance object may be smooth even when a particular spectral coordinate
representation is not.  Classical numerical linear algebra also treats Cholesky
factorization as the stable coordinate chart for strictly positive definite
covariances; see, for example, \citet{higham2002accuracy}.

A practical differentiable variant is therefore to replace the support-exact
eigensystem route by a strictly positive definite, ridged Cholesky gauge.  This
is a different numerical contract, not a proof that the eigensystem contract was
unnecessary.  For the canonical BayesFilter route, the realized scalar ridge
\(\lambda>0\) is selected outside the candidate-dependent differentiated
callable, stored with the prepared inputs, and bound into route identity.  The
same value is used for the value, gradient, finite-difference center, and every
finite-difference endpoint.  Candidate-dependent ridge escalation followed by a
stopped derivative computes a different gradient and is wrong relative to the
canonical total-gradient claim.  A failed Cholesky factor or nonpositive factor
diagonal is an invalid-chart veto, not a trigger for an undisclosed fallback or
escalation.

The canonical route fixes \(\rho=1\).  In row-vector storage, let the prepared,
parameter-independent residual design be
\(\Xi\in\mathbb R^{N\times d_x}\), with \(\mathbf 1^\top\Xi=0\).  Its realized
entries, not merely a seed distribution, are part of the finite target.  For
source rows \(X_i\), normalized probabilities \(w_i\), and row-quotient rows
\(Y_i^+\), define the population moments
\begin{align}
  \mu_w&:=\sum_iw_iX_i,&
  \Sigma_w&:=\sum_iw_i(X_i-\mu_w)(X_i-\mu_w)^\top,\\
  \bar Y^+&:=N^{-1}\sum_iY_i^+,&
  \Sigma_+&:=N^{-1}\sum_i(Y_i^+-\bar Y^+)(Y_i^+-\bar Y^+)^\top,
\end{align}
and \(G_+:=\operatorname{sym}(\Sigma_w-\Sigma_+)\).  Then define
\begin{equation}
\label{eq:bf-eot-contract-e-chol-ridge}
  B_\lambda:=\operatorname{chol}(G_+ + \lambda I),
  \qquad
  \widetilde Y:=Y^+ + \Xi B_\lambda^\top .
\end{equation}
All source and equal-weight covariances in this route are population empirical
covariances with denominator \(N\), not \(N-1\).  The factor
\(\sqrt{N/(N-1)}\), when used to construct \(\Xi\) from centered Gaussian
noise, standardizes that fixed design in expectation and does not change the
covariance convention.
Then form
\[
  \widetilde\Sigma
  :=
  \frac1N
  (\widetilde Y-\mathbf 1\bar{\widetilde y}^\top)^\top
  (\widetilde Y-\mathbf 1\bar{\widetilde y}^\top) ,
  \qquad
  L_w:=\operatorname{chol}(\Sigma_w+\lambda I),
  \qquad
  L_{\widetilde{}}:=\operatorname{chol}(\widetilde\Sigma+\lambda I).
\]
The affine operator is not formed with an inverse explicitly.  It is the
triangular-solve operator
\begin{equation}
\label{eq:bf-eot-contract-e-chol-affine}
  A_\lambda
  :=
  L_w L_{\widetilde{}}^{-1},
  \qquad
  Y^\star_\lambda
  :=
  \mathbf 1\mu_w^\top
  +
  (\widetilde Y-\mathbf 1\bar{\widetilde y}^\top)A_\lambda^\top.
\end{equation}
The recentering in \eqref{eq:bf-eot-contract-e-chol-affine} is essential: it is
what preserves the target mean independently of the ridge.

\begin{proposition}[Ridged Cholesky restoration]
\label{prop:bf-eot-contract-e-chol-restoration}
Assume the Cholesky factors in
\eqref{eq:bf-eot-contract-e-chol-affine} exist.  Then
\[
  A_\lambda(\widetilde\Sigma+\lambda I)A_\lambda^\top
  =
  \Sigma_w+\lambda I,
\]
and \(Y^\star_\lambda\) has equal-weight mean \(\mu_w\).  The covariance of
\(Y^\star_\lambda\) around \(\mu_w\) satisfies
\[
  \frac1N
  (Y^\star_\lambda-\mathbf 1\mu_w^\top)^\top
  (Y^\star_\lambda-\mathbf 1\mu_w^\top)
  -
  \Sigma_w
  =
  \lambda(I-A_\lambda A_\lambda^\top).
\]
\end{proposition}

\begin{proof}
The mean statement follows immediately from the centering in
\eqref{eq:bf-eot-contract-e-chol-affine}.  Since
\(\widetilde\Sigma+\lambda I=L_{\widetilde{}}L_{\widetilde{}}^\top\),
\[
  A_\lambda(\widetilde\Sigma+\lambda I)A_\lambda^\top
  =
  L_wL_{\widetilde{}}^{-1}
  L_{\widetilde{}}L_{\widetilde{}}^\top
  L_{\widetilde{}}^{-\top}L_w^\top
  =
  L_wL_w^\top
  =
  \Sigma_w+\lambda I.
\]
Subtracting
\(A_\lambda(\lambda I)A_\lambda^\top\) from both sides gives the displayed
unridged covariance residual.
\end{proof}

Contract E--Chol is the only reset semantics eligible for canonical BayesFilter
value, total-gradient, score-admission, leaderboard, default, or HMC-facing
status.  This selection does not establish that the route has passed those
scientific gates.  Its advantage is that the derivative can be written through
Cholesky factors and triangular solves, without eigenvector gauges or
rank-switching projectors.  Its cost is the explicit \(\lambda\)-bias above.
Artifacts must report the prepared realized ridge, the exact ridged-identity
residual, and the separate raw-covariance residual.  The proposition proves the
ridged identity; calling the raw covariance exactly restored when
\(\lambda>0\) is wrong.

The positive finite transport feeding this reset must also be interpreted
literally.  If the fixed-iteration streaming solver emits a numerator and row
mass
\[
  Q_i=\sum_j P_{ij}X_j,\qquad M_i=\sum_j P_{ij},
\]
then the transported cloud is
\begin{equation}
\label{eq:bf-eot-contract-e-row-quotient}
  Y_i^+=Q_i/M_i.
\end{equation}
There is no row-mass floor in the canonical target.  Each \(M_i\) must be finite
and positive.  Treating \(Q_i\) as the cloud merely because the finite solver
intends unit row mass changes both the value and its derivative.

At each filtering time the declared transition and finite LEDH flow are applied
first, corrected log weights and the current observed-data likelihood increment
are then computed, and only afterward is an active reset applied.  The reset
affects later increments through the carried cloud; it does not modify the
already-computed current increment.  Active branches carry
\(Y^\star_\lambda\) with uniform log weights.  Inactive branches carry the
flowed cloud with its normalized log weights.  The reset mask, prepared residual
design, ridge, Sinkhorn schedule, and chunk sizes are fixed parts of the finite
program.

\paragraph{Contract F: no reset or delayed reset.}
If the effective sample size is acceptable, carry the weighted cloud forward and
do not apply \(D_B\).  This avoids the deterministic covariance loss at that step.
It does not solve weight degeneracy in long horizons or high dimensions, but it is
a clean diagnostic: if a no-reset weighted route agrees with an exact LGSSM
comparator while the barycentric route fails, the reset contract rather than the
transition or likelihood code is implicated.

\section{Why the current reset can accumulate bias}
\label{sec:bf-eot-contract-e-bias}

Contract~E--Chol is a declared finite filtering program, not an identity for the
exact filtering distribution.  At time \(t\), the flowed particles \(x_{t,i}\)
and previous weights \(w_{t-1,i}\) first form
\begin{equation}
\label{eq:bf-eot-finite-increment}
 r_{t,i}=w_{t-1,i}\,
 \frac{p_\theta(x_{t,i}\mid x_{t-1,i})
       g_\theta(y_t\mid x_{t,i})}
      {q_\theta(x_{t,i}\mid x_{t-1,i})},
 \qquad
 \widehat\ell_t^N=\log\sum_i r_{t,i},
\end{equation}
where the implemented flow Jacobian is included in the ratio.  Only after this
increment has been accumulated does an active reset replace the cloud.  The
reset therefore cannot change the already computed increment, but it can change
all later increments through the carried particles.

There are four separate approximation layers: finite-particle log normalization;
finite entropic Sinkhorn iteration; barycentric projection; and the residual/
affine reset.  Marginal residuals certify only the finite coupling calculation.
They do not remove the barycentric covariance contraction, the finite log-sum
bias, or the loss of higher-order shape information.  Repeated resets can feed
these local errors into later nonlinear transition and observation functionals.
The LGSSM is therefore a controlled diagnostic for this mechanism, not the
scientific endpoint for which the method is intended.

\section{A non-fused cubature residual Contract E candidate}
\label{sec:bf-eot-cubature-residual}

This candidate retains the original staged route.  It does not solve a joint
OT--moment optimization.  The coupling, barycentric quotient, residual stage,
and Cholesky restoration remain separate finite maps; only the prepared
residual design is changed.

Let \(d,M\in\mathbb N\), let \(N=2dM\), and use the population-covariance
convention (denominator \(N\)) throughout this section.  The reset is defined
on a fixed finite-arithmetic branch: the row ordering, the branch of every
Cholesky factor, the ridge \(\lambda>0\), and all Sinkhorn controls are fixed
and independent of the model parameter for the derivative statement below.
The particle-count condition is therefore
\begin{equation}
\label{eq:bf-eot-cubature-particle-count}
 N=2dM.
\end{equation}
Define the standardized spherical--radial cubature directions
\begin{equation}
\label{eq:bf-eot-cubature-directions}
 \xi_{a,+}=\sqrt d\,e_a,\qquad
 \xi_{a,-}=-\sqrt d\,e_a,
 \qquad a=1,\ldots,d,
\end{equation}
and \(D_c=\sqrt d\begin{bmatrix}I_d\\-I_d\end{bmatrix}\).  Let
\(\Xi_c\in\mathbb R^{N\times d}\) repeat every row of \(D_c\) exactly \(M\)
times.  The row order and any fixed independent permutation are part of the
finite route identity.

\begin{proposition}[Cubature design moments]
\label{prop:bf-eot-cubature-design-moments}
For \(\mathbf 1_N\in\mathbb R^N\), the replicated design
\(\Xi_c\in\mathbb R^{N\times d}\) satisfies
\begin{equation}
 \frac1N\mathbf1^\top\Xi_c=0,
 \qquad
 \frac1N\Xi_c^\top\Xi_c=I_d.
\end{equation}
\end{proposition}

\begin{proof}
Opposite directions cancel in every coordinate, so the row sum is zero.  Also
\(D_c^\top D_c=d(I_d+I_d)=2dI_d\).  Repeating the block \(M\) times and using
\(N=2dM\) gives
\[
 N^{-1}\Xi_c^\top\Xi_c=(2dM)^{-1}M(2dI_d)=I_d.
\]
\end{proof}

\begin{proposition}[Affine cubature moment preservation]
\label{prop:bf-eot-cubature-affine-moments}
Let \(a_r=1/N\) and let \(\xi_r\) be the rows of \(\Xi_c\).  For any
  \(m\in\mathbb R^d\), \(L\in\mathbb R^{d\times d}\), \(F\in\mathbb R^{p\times
d}\), and \(b\in\mathbb R^p\), the cubature rule obeys
\begin{align}
 \sum_{r=1}^N a_r\{F(m+L\xi_r)+b\}&=Fm+b,\\
 \sum_{r=1}^N a_r (FL\xi_r)(FL\xi_r)^\top&=FL\,L^\top F^\top.
\end{align}
\end{proposition}

\begin{proof}
The first identity is \(F m+b+FL(N^{-1}\Xi_c^\top\mathbf 1_N)\),
which equals \(Fm+b\) by Proposition~\ref{prop:bf-eot-cubature-design-moments}.
For the second identity, factor out \(FL\) and \(L^\top F^\top\), then use
\(N^{-1}\Xi_c^\top\Xi_c=I_d\).  Thus the rule is exact for the first two
moments of every affine image of a standardized Gaussian moment pair.  This
standalone statement does not imply that the particle log-sum likelihood is
unbiased.
\end{proof}

The OT stage is unchanged.  For flowed source particles \(x_j\) and normalized
probabilities \(\alpha_j\), write the equal-target coupling as \(\Pi\), with
source marginal \(\alpha\) and target marginal \(u_i=1/N\).  The implementation
uses \(\Gamma=N\Pi\), so
\begin{equation}
\label{eq:bf-eot-cubature-marginals}
 \Gamma\mathbf1=\mathbf1,
 \qquad
 \Gamma^\top\mathbf1=N\alpha,
\end{equation}
up to recorded finite-Sinkhorn residuals, and the barycentric output is
\begin{equation}
\label{eq:bf-eot-cubature-barycentric}
 Y_i^+=\sum_j\Gamma_{ij}x_j.
\end{equation}
Cubature is not an OT target in this candidate; it enters only after
\eqref{eq:bf-eot-cubature-barycentric}.

Define
\begin{align}
 \mu_w&=\sum_j\alpha_jx_j,&
 \Sigma_w&=\sum_j\alpha_j(x_j-\mu_w)(x_j-\mu_w)^\top,\\
 \bar Y^+&=N^{-1}\sum_iY_i^+,&
 \Sigma_+&=N^{-1}\sum_i(Y_i^+-\bar Y^+)(Y_i^+-\bar Y^+)^\top.
\end{align}
Set \(G=\operatorname{sym}(\Sigma_w-\Sigma_+)\) and
\(B=\operatorname{chol}(G+\lambda I)\), \(\lambda>0\), then inject
\begin{equation}
\label{eq:bf-eot-cubature-injected}
 \widetilde Y=Y^++\Xi_cB^\top.
\end{equation}
Let \(\widetilde\Sigma\) be its population covariance and define
\begin{equation}
\label{eq:bf-eot-cubature-affine}
 L_w=\operatorname{chol}(\Sigma_w+\lambda I),\quad
 L_{\widetilde Y}=\operatorname{chol}(\widetilde\Sigma+\lambda I),\quad
 A_c=L_wL_{\widetilde Y}^{-1},
\end{equation}
where the inverse denotes a triangular solve.  The reset is
\begin{equation}
 X_i^{E,c}=\mu_w+(\widetilde Y_i-\bar{\widetilde Y})A_c^\top.
\end{equation}

\begin{proposition}[Mean and ridged covariance]
\label{prop:bf-eot-cubature-restoration}
Assume \(G+\lambda I_d\), \(\Sigma_w+\lambda I_d\), and
\(\widetilde\Sigma+\lambda I_d\) are strictly positive definite, so that all
three displayed Cholesky factors exist.  The equal-weight output
\(X^{E,c}\in\mathbb R^{N\times d}\) has mean \(\mu_w\) and satisfies
\begin{equation}
 A_c(\widetilde\Sigma+\lambda I)A_c^\top=\Sigma_w+\lambda I,
 \qquad
 \Sigma_{E,c}-\Sigma_w=\lambda(I-A_cA_c^\top).
\end{equation}
\end{proposition}

\begin{proof}
Because \(\bar{\widetilde Y}=N^{-1}\sum_i\widetilde Y_i\), the centered output
rows sum to zero, and its population covariance is
\(\Sigma_{E,c}=A_c\widetilde\Sigma A_c^\top\).  With
\(L_wL_w^\top=\Sigma_w+\lambda I_d\) and
\(L_{\widetilde Y}L_{\widetilde Y}^\top=\widetilde\Sigma+\lambda I_d\),
\begin{align*}
A_c(\widetilde\Sigma+\lambda I_d)A_c^\top
&=L_wL_{\widetilde Y}^{-1}L_{\widetilde Y}L_{\widetilde Y}^\top
  L_{\widetilde Y}^{-\top}L_w^\top\\
&=L_wL_w^\top=\Sigma_w+\lambda I_d.
\end{align*}
Subtracting \(\lambda A_cA_c^\top\) from both sides gives
\(\Sigma_{E,c}-\Sigma_w=\lambda(I_d-A_cA_c^\top)\).
\end{proof}

The proposition concerns the finite weighted empirical target \(\Sigma_w\), not
the exact filtering covariance.  To make the role of the second stage explicit,
write \(Y_c^+=Y^+-\bar Y^+\) and
\(C_{+c}=N^{-1}(Y_c^+)^\top\Xi_c\).  Proposition
\ref{prop:bf-eot-cubature-design-moments} gives
\begin{equation}
\label{eq:bf-eot-cubature-realized-covariance}
\widetilde\Sigma
 =\Sigma_+ + C_{+c}B^\top + BC_{+c}^\top + BB^\top.
\end{equation}
Thus \(B=\operatorname{chol}(G+\lambda I_d)\) is not by itself a covariance
correction unless \(C_{+c}=0\).  The subsequent \(A_c\) uses the realized
\(\widetilde\Sigma\), including both cross terms, and is the operation that
establishes the ridged identity.

\begin{proposition}[Current likelihood and total-score boundary]
\label{prop:bf-eot-cubature-likelihood-score}
If the cubature residual reset is applied after the current increment
\eqref{eq:bf-eot-finite-increment}, that increment is exactly unchanged relative
to canonical Contract~E.  On a fixed finite-arithmetic branch, with parameter-
independent \(\lambda\) and \(\Xi_c\), differentiating every dependency through
the staged map gives the total derivative of the new cubature-residual finite
likelihood scalar.
\end{proposition}

\begin{proof}
The current increment depends only on pre-reset particles, weights, densities,
and the flow Jacobian.  The full scalar is a composition through
\((x,\alpha,\Pi,Y^+,\mu_w,\Sigma_w,\widetilde\Sigma,A_c,X^{E,c})\).  Since
\(d\Xi_c=d\lambda=0\), the ordinary chain rule through the fixed SPD Cholesky
branches gives the total derivative of every later increment.  This proves a
same-scalar derivative statement for the new finite program, not an
exact-posterior score statement.
\end{proof}

For affine maps \(f(x)=Fx+b\) and \(h(x)=Hx+c\), the preceding affine-moment
proposition gives exact transformed means and covariances for the standalone
cubature rule.  A cubature Gaussian filter therefore agrees with the
corresponding Kalman moment recursion up to arithmetic error when its input is
represented by that rule.  The staged particle route still has finite-particle,
finite-Sinkhorn, barycentric, and reset approximations; this does not make its
particle log-sum estimator unbiased for the Kalman likelihood.

\section{Generalized unscented transformation candidate}
\label{sec:bf-eot-genut}

Ebeigbe et al.~\cite{ebeigbe2025genut} propose a $2d+1$ point generalized
unscented transform (GenUT).  In the source paper, the sigma points match the
mean and covariance exactly and then use the remaining axis freedom to match
the \emph{physical-coordinate diagonal} third and fourth central moments,
constructed through elementwise powers of a chosen covariance square root.
The construction used below keeps the same per-axis algebra but applies it to
whitened moments of $Z=C^{-1}(X-m)$.  For correlated states those whitened
diagonal moments are generally different targets from the paper's physical
diagonal moments, so this section should be read as a GenUT-motivated local
variant rather than as a restatement of the paper's multivariate construction.

Let \(m\in\mathbb R^d\), \(C\in\mathbb R^{d\times d}\) be nonsingular with
\(CC^\top=P\succ0\), and let \(Z=C^{-1}(X-m)\).  Define the coordinate moments
\begin{equation}
 s_a=\mathbb E[Z_a^3],\qquad k_a=\mathbb E[Z_a^4].
\end{equation}
For each axis use standardized points \(0,-u_ae_a,+v_ae_a\), with weights
\(w_0,b_a,c_a\).  The axis constraints are
\begin{align}
 -b_au_a+c_av_a&=0,&
 b_au_a^2+c_av_a^2&=1,\\
 -b_au_a^3+c_av_a^3&=s_a,&
 b_au_a^4+c_av_a^4&=k_a.
\end{align}

\begin{proposition}[GenUT axis construction]
\label{prop:bf-eot-genut-axis}
Fix an axis \(a\).  If \(k_a>s_a^2\), set
\begin{equation}
 u_a=\frac{-s_a+\sqrt{4k_a-3s_a^2}}2,
 \qquad v_a=u_a+s_a,
 \qquad c_a=\frac1{v_a(u_a+v_a)},
 \qquad b_a=\frac1{u_a(u_a+v_a)}.
\end{equation}
Then the four axis constraints hold.  Setting
\(w_0=1-\sum_a(b_a+c_a)\) normalizes the rule, but \(w_0\) is not guaranteed
to be nonnegative.
\end{proposition}

\begin{proof}
The first constraint gives \(b_a=c_av_a/u_a\).  Substitution into the second
gives \(c_av_a(u_a+v_a)=1\).  The third therefore reduces to
\(v_a-u_a=s_a\).  Substitution into the fourth yields
\(u_a^2+s_au_a+s_a^2-k_a=0\), whose positive root is the displayed \(u_a\).
The expression for \(v_a\) and the weights then follow from the first two
constraints.  Since \(\sqrt{4k_a-3s_a^2}>|s_a|\), both \(u_a\) and \(v_a\) are
strictly positive, and therefore \(b_a,c_a>0\).  The central weight supplies
the remaining normalization.  The condition \(k_a>s_a^2\) ensures a real,
positive axis construction, but it says nothing about \(w_0\).
\end{proof}

\begin{proposition}[GenUT moment identities]
\label{prop:bf-eot-genut-moments}
Define \(2d+1\) points \(z_0=0\), \(z_{a,-}=-u_ae_a\), and
\(z_{a,+}=v_ae_a\), with weights \(w_0,b_a,c_a\), and assume
\(w_0=1-\sum_a(b_a+c_a)\) is finite.  Then the GenUT rule has zero mean,
identity covariance, diagonal third moments \(s_a\), and diagonal fourth
moments \(k_a\) in whitened coordinates.  The physical points
\(x_r=m+Cz_r\) therefore have mean \(m\) and covariance \(P\).
\end{proposition}

\begin{proof}
Normalization is the definition of \(w_0\).  For the mean, each axis contributes
\((-b_au_a+c_av_a)e_a=0\).  For the covariance, off-diagonal entries vanish
because every noncentral point has only one nonzero coordinate, while the
\(a\)-th diagonal entry is \(b_au_a^2+c_av_a^2=1\).  The diagonal third and
fourth entries are exactly the last two axis equations; mixed third/fourth
entries involving two distinct coordinates vanish for this rule.  Finally,
\(\sum_r w_r(m+Cz_r)=m+C\sum_r w_rz_r=m\), and
\(\sum_r w_r(x_r-m)(x_r-m)^\top
=C(\sum_r w_rz_rz_r^\top)C^\top=CC^\top=P\).
\end{proof}

\begin{proposition}[GenUT positivity and representation boundary]
\label{prop:bf-eot-genut-positivity}
Under \(k_a>s_a^2\), the noncentral weights \(b_a,c_a\) are positive.  The
full GenUT rule is a positive probability rule if and only if
\(w_0\ge0\).  It can be represented exactly as an equal-weight \(N\)-row
design by replication only when all its positive weights are rational with a
common denominator dividing \(N\) (or when an explicitly declared positive
approximation is used).  A negative \(w_0\) is a signed quadrature rule and
cannot be used as a positive entropic-OT marginal.
\end{proposition}

\begin{proof}
The positivity of \(b_a,c_a\) follows from the strict positivity of
\(u_a,v_a,u_a+v_a\).  Since the weights sum to one, nonnegativity of all
weights is equivalent to \(w_0\ge0\).  Exact equal-weight replication assigns
each row mass \(1/N\), so a weight \(q\) is representable by rows precisely
when \(Nq\) is a nonnegative integer; a common denominator gives the stated
condition.  Signed masses violate the positive-marginal assumptions of
entropic OT.
\end{proof}

For a feasible positive GenUT residual design \(\Xi_g\), obtained by mapping the whitened points \(z_r\) through the declared square root and then using an exact positive equal-weight replication when the condition in Proposition~\ref{prop:bf-eot-genut-positivity} holds, define
\begin{equation}
 \widetilde Y_g=Y^++\Xi_gB_g^\top,
 \qquad B_g=\operatorname{chol}(G_g+\lambda I),
\end{equation}
where the displayed Cholesky argument is assumed strictly positive definite.
Apply the same Cholesky restoration as in
\eqref{eq:bf-eot-cubature-affine}.  This is the GenUT residual candidate.  Its
current likelihood increment is unchanged because the reset is post-increment.
The variance statement is the same ridged identity as in
Proposition~\ref{prop:bf-eot-cubature-restoration}, with \(c\) replaced by
\(g\), provided the realized \(\widetilde\Sigma_g\) is used.  Its total score is
the total derivative of the new finite route when the empirical
third/fourth moments, square-root gauge, point coordinates, weights, and
restoration are all differentiated.  Neither future likelihood equality nor
exact posterior variance/score equality follows.

\begin{proposition}[Value, variance, and score for both staged residuals]
\label{prop:bf-eot-both-residual-boundary}
Let \(q\in\{c,g\}\) denote either the cubature residual or a feasible positive
GenUT residual, and let \(R_q\) be its post-increment reset.  Assume all
Cholesky arguments are strictly positive definite and the parameter is inside
a fixed differentiable branch.  Then:
\begin{enumerate}[label=(\roman*)]
\item the current finite likelihood increment \(\widehat\ell_t^N\) in
\eqref{eq:bf-eot-finite-increment} is identical before and after replacing the
canonical reset by \(R_q\);
\item the reset mean is \(\mu_w\), and its population covariance obeys
\[
\Sigma_{R_q}+\lambda A_qA_q^\top=\Sigma_w+\lambda I_d,
\]
so the un-ridged covariance equals \(\Sigma_w\) only in the special case
\(A_qA_q^\top=I_d\) (or in the zero-ridge limit when that limit exists);
\item differentiating the complete staged computation gives the total
derivative of the complete finite likelihood scalar generated by \(R_q\),
including the derivatives of the OT map, moment maps, residual design,
Cholesky factors, and all later recursion.
\end{enumerate}
The proposition does not assert equality with the exact filtering likelihood,
exact filtering variance, or exact filtering score.  Except in a special
insensitivity case, the full future likelihood scalar for \(R_c\) or \(R_g\) is
different from canonical Contract~E because the carried cloud differs.
\end{proposition}

\begin{proof}
The reset is evaluated after \(\widehat\ell_t^N\), so item (i) is a direct
statement about program order.  Item (ii) is the proof of
Proposition~\ref{prop:bf-eot-cubature-restoration}; the same algebra applies to
any \(q\) whose realized centered design and affine restoration are used.  For
item (iii), on a fixed branch the finite filter is a composition of
differentiable maps.  The chain rule therefore gives the derivative of the
same scalar.  A later increment is a function of the carried reset cloud, so
there is no reason for it to agree with the canonical value unless the later
functional is insensitive to the changed directions.  This is a finite-program
statement, not an exact filtering theorem.
\end{proof}

\begin{center}
\scriptsize
\setlength{\tabcolsep}{3pt}
\begin{tabular}{@{}p{0.20\textwidth} p{0.22\textwidth} p{0.22\textwidth} p{0.24\textwidth}@{}}
\toprule
Property & Cubature residual & Feasible positive GenUT residual & Exactness boundary \\
\midrule
Current likelihood increment & unchanged & unchanged & reset is post-increment; this is a program-order identity \\
Reset mean & \(\mu_w\) & \(\mu_w\) & finite empirical target only \\
Reset variance & same ridged identity & same ridged identity & exact \(\Sigma_w\) only without the ridge residual \\
Higher moments & symmetric second-order design & selected whitened diagonal \(s_a,k_a\) when feasible & neither matches the full correlated filtering law \\
Full future likelihood & generally changed & generally changed & carried cloud changes downstream increments \\
Score & total derivative of cubature finite scalar & total derivative of GenUT finite scalar & not an exact-posterior score claim \\
\bottomrule
\end{tabular}
\setlength{\tabcolsep}{6pt}
\normalsize
\end{center}

The constrained GenUT variant in the paper moves points toward support bounds.
That operation preserves second-order constraints only under its stated
re-solving rules and generally sacrifices exact fourth-moment matching.  A
parameter-dependent constraint slack is therefore a new route control and must
be included in the route identity and derivative graph.

\begin{center}
\scriptsize
\setlength{\tabcolsep}{3pt}
\begin{tabular}{@{}p{0.20\textwidth} p{0.23\textwidth} p{0.22\textwidth} p{0.20\textwidth}@{}}
\toprule
Remedy & What it fixes & Main cost & Required audit \\
\midrule
Same-scalar barycentric contract & claim mismatch for the executed scalar & does not restore covariance & same-scalar finite differences and moment-drift report \\
Stochastic resampling & unbiased empirical-measure semantics in expectation & discontinuous random branch & estimator contract, seed/randomness policy, MCSE \\
Second-order transform & mean and covariance matching & possible negative transform entries and nonphysical states & moment gate, constraint gate, custom VJP gate \\
Affine moment restoration & local covariance restoration diagnostic & new finite target and rank/jitter choices & LGSSM value/gradient gate and state-constraint audit \\
Positive residual-affine hybrid & covariance restoration with positive first stage & novel randomized finite target and conditioning sensitivity & spectrum gate, MCSE, same-scalar gradient gate, constraint audit \\
No-reset threshold & avoids unnecessary deterministic collapse & leaves weight degeneracy unresolved & ESS policy and weighted-cloud comparator \\
\bottomrule
\end{tabular}
\setlength{\tabcolsep}{6pt}
\normalsize
\end{center}

\paragraph{Implication for SIR, SV, and learned teachers.}
Yes, this issue can affect SIR, stochastic-volatility, and other nonlinear-model
gradients.  The effect is harder to expose than in an LGSSM because there is no
closed-form Kalman likelihood oracle, but the mechanism is the same: a
mean-preserving reset can still change nonlinear transition, observation, and
future-likelihood functionals.  Formally, let \(\Phi_t\) denote the rest of the
filtering computation after a reset at time \(t\).  The barycentric route and a
moment-preserving route evaluate different functions,
\[
  L_B(\theta)=\Phi_t\!\left(X_t(\theta)D_B(\theta);\theta\right),
  \qquad
  L_\sharp(\theta)=
  \Phi_t\!\left(X_t(\theta)D^\sharp(\theta);\theta\right),
\]
unless \(\Phi_t\) is insensitive to the covariance directions removed by
\eqref{eq:bf-eot-D-covariance-gap}.  Therefore equality of autodiff and finite
differences for \(L_B\) certifies only the executed barycentric route.  It does
not certify \(L_\sharp\), a categorical-resampling route, or an exact filtering
likelihood.  A learned teacher trained to imitate the current barycentric teacher
\(D_B\) will inherit this moment contract rather than repair it.  A learned or
retained teacher can be a remedy only if its target is changed to one of the
declared objects above, for example \(D^\sharp\), \(D_{\mathrm{NETF}}\), the
stochastic estimator, or the affine-corrected cloud \(Y^\star\), and the
downstream filtering diagnostics are rerun under that new finite target.

\section{Differentiation contracts}
\label{sec:bf-eot-differentiation}

The final major distinction of the chapter is between three different derivative
objects:
\begin{enumerate}[label=(\roman*)]
  \item derivative of the exact regularized OT optimizer,
  \item derivative of the finite unrolled Sinkhorn computation,
  \item derivative of some scalar built from that finite computation.
\end{enumerate}
In code, the second and third are the most common.  A gradient is only the
correct gradient of the exact executed scalar if the computational graph used in
autodiff matches the scalar value actually reported by the code.  The next chapter
specializes this point to the declared LEDH-PFPF-OT scalar and its custom VJP.

\begin{proposition}[Stopped transport normalization gives a partial derivative]
\label{prop:bf-eot-stopped-normalization-partial}
Let a finite Sinkhorn transport step use normalized particles
\[
  u_i(x)=\frac{x_i-\bar x}{s(x)},\qquad
  \bar x=\frac{1}{N}\sum_{k=1}^N x_k,
\]
where \(s(x)\) is the scale used to build the transport cost.  Let
\[
  L(x)=\Psi\bigl(x,u(x)\bigr)
\]
be any downstream scalar depending on both the carried particles \(x\) and the
transport-normalized particles \(u(x)\).  If a reverse pass treats \(\bar x\)
and \(s(x)\) as constants, then it computes the partial derivative of
\(\Psi(x,u)\) with respect to \(x\) under the artificial map
\[
  u_i=\frac{x_i-c}{s},
\]
not the total derivative of \(L(x)\).  Therefore, unless the omitted terms are
zero or a different objective is explicitly declared, this stopped derivative is
wrong relative to the claim that it is the score of the executed finite scalar.
\end{proposition}

\begin{proof}
Write \(a_i=\partial L/\partial u_i\).  The true differential through
\[
  u_i(x)=\frac{x_i-\bar x}{s(x)}
\]
is
\[
  dL_u
  =
  \sum_i a_i^\top \frac{dx_i-d\bar x}{s(x)}
  -
  \left(\sum_i a_i^\top (x_i-\bar x)\right)\frac{ds(x)}{s(x)^2}.
\]
Since \(d\bar x=N^{-1}\sum_k dx_k\), the contribution to the adjoint of
\(x_i\) is
\[
  \frac{a_i}{s(x)}
  -
  \frac{1}{N}\sum_k \frac{a_k}{s(x)}
  -
  \frac{1}{s(x)^2}
  \left(\sum_k a_k^\top (x_k-\bar x)\right)\nabla_{x_i}s(x).
\]
The stopped rule keeps only \(a_i/s(x)\).  It drops the mean term and the scale
term.  Hence it is the derivative of an artificial partial-derivative problem
with fixed \(c\) and fixed \(s\), not the total derivative of \(L(x)\).
\end{proof}

The same warning applies to Sinkhorn cost keys, adaptive starting
\(\varepsilon_0\), finite-iteration potentials, and any other quantity computed
from the current particles.  If the scalar value recomputes such a quantity
when \(\theta\) changes, then its score includes the corresponding chain-rule
term.  A hand-written VJP may use internal stops only to avoid double counting,
but it must add back every blocked derivative.  Otherwise the result is a
partial derivative and must be named as such.

\begin{center}
\small
\begin{tabular}{p{0.24\textwidth} p{0.30\textwidth} p{0.33\textwidth}}
\toprule
Derivative object & What is being differentiated? & What can go wrong? \\
\midrule
Exact regularized optimizer & the ideal entropic OT solution & not what finite code usually executes directly \\
Finite unrolled Sinkhorn route & the declared numerical balancing procedure & residuals or stabilization may change the executed object \\
Downstream scalar built from finite Sinkhorn & the actual reported value used in filtering or HMC-adjacent code & value-gradient mismatch if autodiff follows a different graph \\
\bottomrule
\end{tabular}
\end{center}

\section{Filtering-loop insertion point}
\label{sec:bf-eot-loop}

Inside a filter, the finite EOT/Sinkhorn equal-weighting step sits after the
weighted particle cloud has been formed and before an equally weighted cloud is
passed onward:
\begin{equation}
\label{eq:bf-eot-loop}
  \{x_i,w_i\}_{i=1}^N
  \xrightarrow{\mathrm{EOT/Sinkhorn}}
  \{\tilde x^{(j)},1/N\}_{j=1}^N.
\end{equation}
If the broader filtering algorithm carries auxiliary particle-local state, this
chapter does not yet assert how that state is transported.  Its contract is
narrower: it specifies the equal-weighting rule for the particle locations and the
mass-routing matrix from which that equal-weight cloud is derived.

\section{Verification contract}
\label{sec:bf-eot-verification}

A reader implementing this route should verify at least:
\begin{enumerate}[label=(\arabic*)]
  \item row and column residuals are reported and decrease with additional
  iterations or stronger stabilization,
  \item the initialization, \(\varepsilon\), iteration budget, and stabilization
  policy are recorded,
  \item the barycentric cloud has the declared shape and normalization,
  \item any reported gradient is tied to the exact computational graph that was
  actually executed,
  \item any stronger filtering or HMC claim is kept separate from these local
  numerical checks.
\end{enumerate}

A failed residual check is evidence about the finite solver or stabilization path,
not about the abstract EOT object.  A successful residual check is evidence that
the numerical route is behaving as declared, not that it preserves the original
categorical target.

\section{What this chapter does and does not claim}
\label{sec:bf-eot-boundary}

This chapter teaches the main differentiable OT computational route in a near-
implementation style.  The derivation above fixes a specific ladder:
\begin{enumerate}[label=(\roman*)]
  \item deterministic OT chooses an exact coupling for the unregularized baseline,
  \item entropic OT changes that teacher object,
  \item Sinkhorn balancing computes a finite numerical approximation to the entropic
    teacher,
  \item barycentric projection turns that finite coupling into the equal-weight
    cloud consumed by later filtering code.
\end{enumerate}
Within that ladder, exact-online and streamed execution strategies remain on the
reference lane when they preserve the same scaling equations and the same
barycentric output contract.  Later retained-teacher chapters ask whether the same
teacher can be compressed or approximated while still judged against this reference
object.  Still later changed-object chapters ask a different question again: what
happens when the transport object itself is replaced by a learned map, altered
support, or new target contract.

It does \emph{not} claim:
\begin{itemize}
  \item that entropic OT is the same object as unregularized OT,
  \item that finite Sinkhorn is the exact regularized optimizer,
  \item that a barycentric equal-weight cloud is the same object as categorical
  ancestor sampling,
  \item that autodiff through the executed graph is automatically the gradient of
    the original filtering likelihood.
\end{itemize}
Its role is to fix the entropic teacher object, the finite numerical route, the
output cloud construction, and the differentiation contract that the recovered
custom-gradient bridge and later learned-transport chapters build on.

\FloatBarrier

\section{A complete higher-moment Contract E candidate}
\label{sec:bf-eot-higher-moment-contract-e}

This section defines the complete opt-in algorithm that is implemented and
tested below.  It is important to name the object correctly.  A classical
unscented transform replaces the distribution of an input \(x\) by sigma
points and evaluates the nonlinear map \(f\) at those sigma points.  The
algorithm here instead propagates the particle rows through \(f\), evaluates
their likelihoods, applies entropic OT and Contract~E, and then applies a
finite higher-moment correction to the carried cloud.  GenUT supplies the motivation and the axis algebra, but this is not the
classical GenUT quadrature algorithm, and the whitened diagonal targets used
here are not the same multivariate moment targets stated in the source paper.  The existing Contract~E route remains
unchanged when correction is disabled.  The optional protected no-fit
comparator below additionally applies its declared cap and affine restoration.

\subsection{Finite filtering, OT, and Contract E stages}

Let \(x_{t-1,i}^{+}\in\mathbb R^d\), \(i=1,\ldots,N\), be the equally
weighted cloud carried from time \(t-1\).  The innovation rows \(e_{t,i}\),
their ordering, and all numerical controls are fixed as part of the finite
route.  The transition and likelihood stages are
\begin{align}
 x_{t,i}^{-}&=f_\theta(x_{t-1,i}^{+},e_{t,i}),
 \label{eq:bf-eot-hm-transition}\\
 \ell_{t,i}&=\log g_\theta(y_t\mid x_{t,i}^{-}),
 &\widehat\ell_t^N&=\log\!\left(\frac1N\sum_{i=1}^N
 e^{\ell_{t,i}}\right),
 \label{eq:bf-eot-hm-likelihood}\\
 \alpha_{t,i}&=
 \frac{e^{\ell_{t,i}}}{\sum_{j=1}^N e^{\ell_{t,j}}}.
 \label{eq:bf-eot-hm-weights}
\end{align}
The weighted empirical target, its mean, and its population covariance are
\begin{align}
 \widehat\pi_t^w&=\sum_{i=1}^N\alpha_{t,i}\delta_{x_{t,i}^{-}},
 &\mu_w&=\sum_{i=1}^N\alpha_{t,i}x_{t,i}^{-},\\
 P_w&=\sum_{i=1}^N\alpha_{t,i}
 (x_{t,i}^{-}-\mu_w)(x_{t,i}^{-}-\mu_w)^\top.
 \label{eq:bf-eot-hm-weighted-moments}
\end{align}

For completeness, the implemented finite Sinkhorn stage uses the cost
\(C_{ij}=\lVert x_{t,i}^{-}-x_{t,j}^{-}\rVert_2^2\), a positive kernel
\(K_{ij}=\exp\{-C_{ij}/(\epsilon\bar C)\}\), where
\(\bar C=\max(N^{-2}\sum_{ij}C_{ij},10^{-3})\), and fixed iteration counts.
Writing \(u_i=1/N\), each iteration applies
\begin{equation}
 a\leftarrow u\oslash(Kb+\tau\mathbf1),
 \qquad
 b\leftarrow\alpha\oslash(K^\top a+\tau\mathbf1),
 \qquad \tau=10^{-7},
 \label{eq:bf-eot-hm-sinkhorn}
\end{equation}
and returns the finite nonnegative coupling
\(\Pi=\operatorname{diag}(a)K\operatorname{diag}(b)\).  Because a finite
iteration need not satisfy either marginal exactly, the carried barycenter is
the row quotient
\begin{equation}
 m_i=\sum_j\Pi_{ij},
 \qquad
 y_i^+=\frac{\sum_j\Pi_{ij}x_{t,j}^{-}}{m_i}.
 \label{eq:bf-eot-hm-barycentric}
\end{equation}
The route fails closed unless every \(m_i>\tau\) and the source marginal of
the quotient coupling
\(\overline\Pi_{ij}=u_i\Pi_{ij}/m_i\) satisfies its declared tolerance.
Replacing the quotient by \(N\sum_j\Pi_{ij}x_j\) is valid only when the finite
row marginal is already exactly \(u\).

Let \(\bar y=N^{-1}\sum_i y_i^+\) and
\(P_+=N^{-1}\sum_i(y_i^+-\bar y)(y_i^+-\bar y)^\top\).  Contract~E uses a
fixed zero-mean residual design \(\Xi\in\mathbb R^{N\times d}\) and
\begin{align}
 G&=\operatorname{sym}(P_w-P_+),
 &B&=\operatorname{chol}(G+\lambda I_d),\\
 \widetilde y_i&=y_i^++\Xi_iB^\top,
 &\overline{\widetilde y}&=N^{-1}\sum_i\widetilde y_i.
 \label{eq:bf-eot-hm-residual}
\end{align}
If \(P_\Xi=N^{-1}\Xi^\top\Xi\) and
\(C_{y\Xi}=N^{-1}\sum_i(y_i^+-\bar y)\Xi_i^\top\), then its realized
population covariance is
\begin{equation}
 \widetilde P=P_++C_{y\Xi}B^\top+BC_{y\Xi}^\top+BP_\Xi B^\top.
 \label{eq:bf-eot-hm-cross-covariance}
\end{equation}
Thus \(BB^\top\) alone is not the injected covariance when the cross term is
nonzero.  The second Contract~E stage uses the realized \(\widetilde P\):
\begin{align}
 L_wL_w^\top&=P_w+\lambda I_d,
 &L_EL_E^\top&=\widetilde P+\lambda I_d,\\
 A&=L_wL_E^{-1},
 &x_i^E&=\mu_w+(\widetilde y_i-\overline{\widetilde y})A^\top.
 \label{eq:bf-eot-hm-contract-e}
\end{align}
Every inverse in this section denotes a triangular solve.

\begin{proposition}[Contract E mean and ridged covariance]
\label{prop:bf-eot-hm-contract-e}
Suppose \(G+\lambda I_d\), \(P_w+\lambda I_d\), and
\(\widetilde P+\lambda I_d\) are strictly positive definite.  Under the
population-covariance convention, the cloud in
\eqref{eq:bf-eot-hm-contract-e} has mean \(\mu_w\) and covariance \(P_E\)
satisfying
\begin{equation}
 P_E+\lambda AA^\top=P_w+\lambda I_d,
 \qquad
 P_E-P_w=\lambda(I_d-AA^\top).
 \label{eq:bf-eot-hm-ridged-identity}
\end{equation}
\end{proposition}

\begin{proof}
The centered rows in \eqref{eq:bf-eot-hm-contract-e} sum to zero, so their
affine image has mean \(\mu_w\).  Its covariance is
\(P_E=A\widetilde P A^\top\).  The Cholesky definitions give
\[
 A(\widetilde P+\lambda I_d)A^\top
 =L_wL_E^{-1}L_EL_E^\top L_E^{-\top}L_w^\top
 =L_wL_w^\top=P_w+\lambda I_d.
\]
Subtraction of \(\lambda AA^\top\) proves both identities.  In particular,
the first stage alone does not prove covariance restoration, and positive
ridge does not generally give the un-ridged equality \(P_E=P_w\).
\end{proof}

\begin{proposition}[Current-increment value boundary]
\label{prop:bf-eot-hm-current-value}
Any reset applied after \eqref{eq:bf-eot-hm-likelihood} leaves the current
increment \(\widehat\ell_t^N\) exactly unchanged.  It generally changes later
increments and therefore generally changes the complete finite filtering
value.
\end{proposition}

\begin{proof}
The current increment is a function only of the pre-reset rows
\(x_{t,i}^{-}\) and their likelihoods.  A reset changes the rows supplied to
the next transition.  Unless all later transition and likelihood functionals
are insensitive to that change, a later increment changes.  This is a
program-order identity for the finite estimator, not equality with the exact
filtering likelihood.
\end{proof}

\subsection{Finite higher-moment correction}

The construction below is deliberately narrower than two established
higher-order sigma-point rules.  Ebeigbe et al.'s GenUT uses \(2d+1\) weighted
points to match the mean, covariance, and selected diagonal third and fourth
moments when its moment and positivity conditions hold
\citep[Secs.~III--V, Algorithms~1--2]{ebeigbe2025genut}.  Under state
constraints, their Algorithm~2 preserves the lower-order construction but
explicitly accepts that exact diagonal kurtosis matching can be lost.  The
higher-order unscented transform (HOUT) of
\citet[Definition~4.1, Theorem~4.2, Algorithm~4.1]{easley2021higherorder}
instead decomposes the full third- and fourth-order tensors into rank-one terms
and builds a variable-size signed-weight quadrature rule.  It has controlled
moment error, but Remark~4.4 shows that its absolute quadrature condition
number can grow like the reciprocal requested tolerance.  Neither method is a
drop-in positive equal-weight map of the existing \(N\)-particle Contract~E
cloud.  We therefore retain the weighted source moments as a least-squares
target and borrow the established principles that constraints can preclude
exact higher-moment matching and that the numerical correction must be damped.
The exact finite composition below is a BayesFilter candidate, not a theorem
claimed by either source.

Assume now that the un-ridged weighted covariance is strictly positive
definite and write \(C_wC_w^\top=P_w\).  With row-vector notation, define
\begin{equation}
 z_i=(x_{t,i}^{-}-\mu_w)C_w^{-\top},
 \qquad
 s_a=\sum_i\alpha_{t,i}z_{ia}^3,
 \qquad
 k_a=\sum_i\alpha_{t,i}z_{ia}^4.
 \label{eq:bf-eot-hm-target-moments}
\end{equation}
These are the selected standardized diagonal moments motivated by the GenUT
axis rule of Proposition~\ref{prop:bf-eot-genut-axis}.  Full third- and
fourth-order tensors are not formed.  Selected mixed entries are retained by
the pairwise correction in Section~\ref{sec:bf-eot-hm-pairwise-repair}.

Before attempting a correction, the implementation records necessary
one-dimensional realizability diagnostics.  For any standardized scalar
random variable, Pearson's moment inequality gives
\begin{equation}
 k_a\geq s_a^2+1.
 \label{eq:bf-eot-hm-pearson-bound}
\end{equation}
For an equal-weight standardized cloud of exactly \(N>1\) rows, a convenient
self-contained necessary upper bound is
\begin{equation}
 k_a\leq N-1.
 \label{eq:bf-eot-hm-finite-n-kurtosis-bound}
\end{equation}
Bounds on the corresponding sample skewness and kurtosis are classical
\citep{johnson1979bounds}; that source gives sharper joint bounds, but the
closed primary text was not available for technical inspection in this audit,
so the implementation uses only the following directly proved inequality.  A
short derivation of
\eqref{eq:bf-eot-hm-finite-n-kurtosis-bound} uses the zero-sum constraint.
Fix \(u_1=x\).  Among the remaining \(N-1\) rows with sum \(-x\),
Cauchy--Schwarz implies
\(
 \sum_{i=2}^N u_i^2\geq x^2/(N-1)
\).
Since \(\sum_i u_i^2=N\), this gives
\(
 x^2\leq N-1
\).  The same argument applies to every row, hence
\(u_i^2\leq N-1\) and
\[
 \frac1N\sum_i u_i^4
 \leq \frac1N (N-1)\sum_i u_i^2
 =N-1.
\]
This bound is conservative rather than sharp, which is appropriate for a
necessary-feasibility diagnostic.

Equations~\eqref{eq:bf-eot-hm-pearson-bound} and
\eqref{eq:bf-eot-hm-finite-n-kurtosis-bound} are necessary marginal conditions,
not sufficient conditions for a jointly realizable \(d\)-dimensional cloud.
The repair does not clip \(s_a\) or \(k_a\), because doing so would silently
change the empirical moment teacher and hence the finite value program.  An
infeasible target remains the declared least-squares target, but exact matching
is forbidden as a claim and every finite displacement is bounded below.

At one correction iteration, standardize the current equal-weight rows to
\(u_{ia}\) so that \(N^{-1}\sum_i u_i=0\) and
\(N^{-1}\sum_i u_i u_i^\top=I_d\).  Let
\begin{equation}
 m_{r,a}=\frac1N\sum_i u_{ia}^r,
 \qquad d_a=\begin{bmatrix}s_a-m_{3,a}\\ k_a-m_{4,a}\end{bmatrix}.
 \label{eq:bf-eot-hm-moment-residual}
\end{equation}
For each coordinate define two projected moment directions
\begin{align}
 q_{3,ia}&=u_{ia}^2-1-m_{3,a}u_{ia},\\
 q_{4,ia}&=u_{ia}^3-m_{3,a}-m_{4,a}u_{ia}.
 \label{eq:bf-eot-hm-directions}
\end{align}
The \(2\times2\) local moment Jacobian is
\begin{equation}
 J_a=
 \begin{bmatrix}
 3N^{-1}\sum_i u_{ia}^2q_{3,ia} &
 3N^{-1}\sum_i u_{ia}^2q_{4,ia}\\
 4N^{-1}\sum_i u_{ia}^3q_{3,ia} &
 4N^{-1}\sum_i u_{ia}^3q_{4,ia}
 \end{bmatrix}.
 \label{eq:bf-eot-hm-jacobian}
\end{equation}
The historical update used unscaled normal equations with an absolute ridge.
That is unsafe when one column of \(J_a\) is many orders of magnitude larger
than another: forming \(J_a^\top J_a\) squares the scale disparity and a fixed
ridge can become numerically irrelevant.  Define instead
\begin{align}
 h_{a,j}&=\sqrt{\lVert J_{a,:,j}\rVert_2^2+\varepsilon_s^2},
 &H_a&=\operatorname{diag}(h_{a,1},h_{a,2}),\\
 A_a&=J_aH_a^{-1},
 &M_a&=A_a^\top A_a+\lambda I_2,
 \label{eq:bf-eot-hm-scaled-lm-system}
\end{align}
where \(\varepsilon_s>0\) is a column-scale floor and \(\lambda>0\) is a
relative Levenberg--Marquardt damping parameter.  The scaled finite coefficient
is
\begin{align}
 y_a&=M_a^{-1}A_a^\top d_a,
 &c_a&=\rho H_a^{-1}y_a,
 \label{eq:bf-eot-hm-coefficients}\\
 \Delta_{ia}&=c_{a,1}q_{3,ia}+c_{a,2}q_{4,ia}.
 \label{eq:bf-eot-hm-correction}
\end{align}
This is a column-scaled Levenberg--Marquardt step for the local nonlinear
least-squares problem \citep{marquardt1963algorithm}.  Column scaling is not
an ad hoc addition: \citet[Sec.~3, steps (iii) and Note (ii)]{osborne1976levenberg}
recommends scaling the Jacobian columns to unit length so that the terms added
to the normal matrix are comparable to its original elements.  If
\(\varepsilon_s=0\) and both columns are nonzero, the columns of \(A_a\) have
unit norm.  Therefore \(\operatorname{tr}(A_a^\top A_a)=2\), its eigenvalues
lie in \([0,2]\), and
\begin{equation}
 \kappa_2(M_a)\leq\frac{2+\lambda}{\lambda}.
 \label{eq:bf-eot-hm-scaled-lm-condition-bound}
\end{equation}
With a positive scale floor the column norms are at most one, so the same upper
bound remains valid.  This bound concerns the scaled two-by-two solve; it does
not prove the nonlinear moment problem is feasible.

To bound the actual cloud move, let the RMS standardized displacement of row
\(i\) be
\begin{equation}
 r_i=\sqrt{d^{-1}\sum_a\Delta_{ia}^2},
 \qquad
 g_i=\left(1+\frac{r_i^2}{\tau^2}\right)^{-1/2},
 \qquad
 v_{ia}=u_{ia}+g_i\Delta_{ia},
 \label{eq:bf-eot-hm-smooth-trust-cap}
\end{equation}
for trust radius \(\tau>0\).  The map is smooth for finite inputs and its
post-cap RMS is \(r_i/\sqrt{1+r_i^2/\tau^2}<\tau\).  Thus even an infeasible
residual or a large local coefficient cannot move any row farther than the
declared radius before rewhitening.  This is not a bound on the complete
iteration after its covariance rescaling.  The corrected \(v_i\) are re-centered and
re-whitened before the next finite iteration.  After the declared number \(S\)
of iterations, the standardized rows \(u_i^H\) are mapped back by
\begin{equation}
 x_i^H=\mu_w+u_i^H C_w^\top.
 \label{eq:bf-eot-hm-map-back}
\end{equation}

\begin{proposition}[Projected directions and local moment Jacobian]
\label{prop:bf-eot-hm-projected-jacobian}
For each coordinate, both directions in
\eqref{eq:bf-eot-hm-directions} have zero empirical mean and zero empirical
inner product with \(u_a\).  Moreover, \(J_a\) in
\eqref{eq:bf-eot-hm-jacobian} is the derivative at \(c_a=0\) of the map from
the two correction coefficients to the third and fourth empirical moments
before rewhitening.
\end{proposition}

\begin{proof}
Because \(m_{1,a}=0\) and \(m_{2,a}=1\),
\[
 N^{-1}\sum_iq_{3,ia}=m_{2,a}-1-m_{3,a}m_{1,a}=0,
 \quad
 N^{-1}\sum_i u_{ia}q_{3,ia}=m_{3,a}-m_{1,a}-m_{3,a}m_{2,a}=0.
\]
The corresponding expressions for \(q_4\) are
\(m_{3,a}-m_{3,a}-m_{4,a}m_{1,a}=0\) and
\(m_{4,a}-m_{3,a}m_{1,a}-m_{4,a}m_{2,a}=0\).
For \(v=u+c_1q_3+c_2q_4\), differentiation of
\(N^{-1}\sum_i v_i^3\) and \(N^{-1}\sum_i v_i^4\) at \(c=0\) gives
respectively \(3N^{-1}\sum_i u_i^2q_j\) and
\(4N^{-1}\sum_i u_i^3q_j\), which are exactly the rows of \(J_a\).
Thus each direction preserves its coordinate mean and variance to first
order; cross-coordinate covariances need not be preserved by the marginal
proposal.  Un-ridged full-rank rewhitening restores the full covariance
exactly.  The numerical relative-ridge convention is distinguished below.
\end{proof}

\begin{proposition}[Exact first and second moments after rewhitening]
\label{prop:bf-eot-hm-rewhitening}
Let \(v_i\in\mathbb R^d\) have population mean \(\bar v\), strictly positive
definite population covariance \(Q\), and Cholesky factor \(RR^\top=Q\).
Set \(u_i=(v_i-\bar v)R^{-\top}\).  Then
\begin{equation}
 \frac1N\sum_i u_i=0,
 \qquad
 \frac1N\sum_i u_i u_i^\top=I_d.
\end{equation}
Consequently \eqref{eq:bf-eot-hm-map-back} has mean \(\mu_w\) and exactly
the un-ridged population covariance \(P_w\).
\end{proposition}

\begin{proof}
Centering proves the first identity.  Row-vector multiplication gives
\[
 \frac1N\sum_i u_i^\top u_i=R^{-1}QR^{-\top}=I_d.
\]
The affine map in \eqref{eq:bf-eot-hm-map-back} therefore has covariance
\(C_wI_dC_w^\top=P_w\), and its centered rows sum to zero.
\end{proof}

The update \eqref{eq:bf-eot-hm-coefficients} is a scaled damped local
Gauss--Newton/Levenberg--Marquardt step, not an exact projection.  The smooth
cap \eqref{eq:bf-eot-hm-smooth-trust-cap} bounds the executed move but can
reduce moment progress.  Rewhitening can change the third and fourth moments,
and the nonlinear moment-feasibility problem can be infeasible or have multiple
solutions.  The implementation therefore reports the necessary-feasibility
diagnostics, pre/post-cap displacement, and post-rewhitening residuals, and
makes no exact higher-moment matching claim.

\subsection{Why the residual design can obstruct kurtosis repair}
\label{sec:bf-eot-hm-richer-design}

A correction cannot change a moment unless its particle configuration supplies
a direction in which that moment can move.  This is a concrete obstruction for
the scalar replicated-axis residual design.  When the OT barycenters have
nearly collapsed, residual injection dominates, and standardization of repeated
\(-1,+1\) residuals gives a cloud close to two symmetric points.  For the exact
binary cloud,
\[
 \overline u=0,\qquad \overline{u^2}=1,\qquad
 m_3=0,\qquad m_4=1.
\]
Substitution in \eqref{eq:bf-eot-hm-directions} gives
\(q_3=u^2-1=0\) and \(q_4=u^3-u=0\) at every particle.
Consequently \(J=0\), and a damped solve returns a zero displacement even when
the target kurtosis is near three.  A larger correction strength cannot repair
this exact degeneracy.  The obstruction concerns this design and these
correction directions; it does not say that every cloud produced by binary
residual injection must itself be exactly binary.

For any zero-mean, unit-variance scalar cloud there is also the useful identity
\begin{equation}
 m_4-1=\operatorname{Var}_{N}(u^2).
 \label{eq:bf-eot-hm-kurtosis-radii}
\end{equation}
Indeed, \(\operatorname{Var}_{N}(u^2)=\overline{u^4}-(\overline{u^2})^2\).
Kurtosis measures variation in squared distance from the mean.  Binary points
all have the same distance and therefore attain the lower value one.  To
represent larger kurtosis at the same variance, the cloud needs different
radii: many smaller distances and some larger ones.  Replacing the binary
residual design supplies those radii before the moment fit begins.

The richer option uses a \emph{fixed finite normal-quantile design}, rather
than sampling a new continuous random variable at each reset.  In one dimension,
let \(\Phi\) be the standard normal distribution function and set
\begin{equation}
 b_i=\Phi^{-1}\!\left(\frac{i+1/2}{N}\right),\quad i=0,\ldots,N-1,
 \qquad
 \Xi_i=\frac{b_i-\bar b}{\sqrt{N^{-1}\sum_j(b_j-\bar b)^2}}.
 \label{eq:bf-eot-hm-normal-quantile-design}
\end{equation}
The implementation arranges opposite-sign pairs.  In several dimensions it
uses fixed permutations of the quantile coordinates, then centers and whitens
the entire design so that \(N^{-1}\Xi^\top\Xi=I_d\), provided its covariance
has full rank.  Both the ordering and the design are frozen for a finite
value/score evaluation, so \(\mathrm d\Xi/\mathrm d\theta=0\).
The implemented construction retains the divisibility requirement \(2d\mid N\).
It is still an empirical distribution with \(N\) atoms; ``richer'' describes
its range of coordinate magnitudes.  In particular, finite quantiles do not
have exactly Gaussian fourth moments, and subsequent residual injection need
not preserve their moments because the barycenters and cross covariance in
\eqref{eq:bf-eot-hm-cross-covariance} also contribute.

A September 2026 scalar KSC diagnostic separated this design change from the
coordinate-cap change described below.  The weighted input kurtosis at the
examined reset was 2.916221.  The same reset, under four choices, gave:
\begin{center}
\begin{tabular}{lrrr}
\toprule
Residual design and coordinate cap & Raw reset & After fit & Final reset\\
\midrule
Repeated axes, original cap & 1.000473 & 1.577655 & 1.328293\\
Normal quantiles, original cap & 2.972274 & 2.966277 & 1.586420\\
Repeated axes, identity-core cap & 1.000473 & 1.577655 & 1.577655\\
Normal quantiles, identity-core cap & 2.972274 & 2.966277 & 2.966277\\
\bottomrule
\end{tabular}
\end{center}
These are standardized fourth moments at one saved reset, not estimates of a
population-wide advantage.  They expose two different losses of shape: the
binary residual design starts near kurtosis one, while the original cap can
erase the spread supplied by the richer design.  The saved reset values are recorded in the September~29 KSC
repair note.\footnote{\path{docs/benchmarks/sqmc-ksc-reset-repair-results-20260929.md}}
The \(T=120\) comparison gives a more informative conclusion than its
binary accuracy screen: the revised inverse-CDF route reduced score error on
both examined datasets, while likelihood error fell on one and rose on the
other.  Section~\ref{sec:bf-eot-hm-observed-results} gives the actual values and
their uncertainty.  A repaired reset moment and a repaired filtering likelihood
remain different empirical questions.

\subsection{The pairwise skewness and kurtosis correction remains active}
\label{sec:bf-eot-hm-pairwise-repair}

Marginal moments cannot describe all dependence between state coordinates.
For the whitened weighted source \(z_i\) in
\eqref{eq:bf-eot-hm-target-moments}, the pairwise correction retains the targets
\begin{equation}
 S^\star_{ab}=\sum_i\alpha_i z_{ia}^{2}z_{ib},\qquad
 K^\star_{ab}=\sum_i\alpha_i z_{ia}^{2}z_{ib}^{2},\qquad a\ne b.
 \label{eq:bf-eot-hm-pair-targets}
\end{equation}
Co-skewness \(S^\star_{ab}\) is ordered: generally
\(S^\star_{ab}\ne S^\star_{ba}\).  Co-kurtosis \(K^\star\) is symmetric.
These selected mixed moments supplement the marginal moments; they are not the
complete third- and fourth-order tensors.  For example,
\(\mathbb E[z_a^3z_b]\) and entries involving three or four distinct coordinates
are outside this pairwise fit.  An explicitly supplied moment teacher may
replace the weighted empirical targets, but that changes what is being fitted
and needs its own accuracy evidence.

Let \(U\) be the current \(N\times d\) standardized cloud, let
\(U^{\circ2}\) denote entrywise squares, and let \(M_3,M_4\) be fixed selection
masks, with zero diagonal.  The all-pairs option sets every off-diagonal entry
to one.  Define
\begin{align}
 S(U)&=N^{-1}(U^{\circ2})^\top U,&
 K(U)&=N^{-1}(U^{\circ2})^\top U^{\circ2},\\
 R_3&=M_3\odot(S^\star-S(U)),&
 R_4&=M_4\odot(K^\star-K(U)).
 \label{eq:bf-eot-hm-pair-residuals}
\end{align}
The implemented direction before projection is
\begin{equation}
 G=\frac{2U\odot(UR_3^\top)+(U^{\circ2})R_3
             +2U\odot((U^{\circ2})R_4^\top)}{\max(d-1,1)}.
 \label{eq:bf-eot-hm-pair-direction}
\end{equation}
For binary masks and symmetric \(M_4\), this equals
\(-N/\max(d-1,1)\) times the gradient with respect to \(U\) of
\(\tfrac12\lVert R_3\rVert_F^2+\tfrac14\lVert R_4\rVert_F^2\).
The factor one quarter accounts for the two occurrences of a symmetric
co-kurtosis entry.  To check the factors directly,
\(\partial S_{ab}/\partial u_{ic}
=N^{-1}(2u_{ia}u_{ib}\mathbf1_{c=a}+u_{ia}^2\mathbf1_{c=b})\);
co-kurtosis has the corresponding two squared-coordinate terms.  Summing
these derivatives gives \eqref{eq:bf-eot-hm-pair-direction}.

The next projection removes the mean and the first-order covariance change:
\begin{equation}
 G_c=G-\mathbf1\bar G,\qquad
 A=\operatorname{sym}(N^{-1}U^\top G_c),\qquad H=G_c-UA.
 \label{eq:bf-eot-hm-pair-projection}
\end{equation}
If \(\bar U=0\) and \(N^{-1}U^\top U=I\), then \(\bar H=0\) and
\(N^{-1}(U^\top H+H^\top U)=0\): substituting \(H\) cancels
\(2A\).  The projection therefore preserves mean and covariance to first
order, although a finite step still needs rewhitening.  With a positive floor
\(\varepsilon_p\), normalize and cap the direction as
\begin{equation}
 V=\frac{H}{\sqrt{(Nd)^{-1}\lVert H\rVert_F^2+\varepsilon_p}},\quad
 r_i^2=d^{-1}\lVert V_i\rVert_2^2,\quad
 D_i=\eta_p\frac{V_i}{\sqrt{1+r_i^2/c_p^2}}.
 \label{eq:bf-eot-hm-pair-cap}
\end{equation}
Each proposed displacement has RMS less than \(\eta_p c_p\).  The candidate
\(U+D\) is then whitened.  Thus the existing pairwise cap remains a
\emph{displacement} cap; changing the final coordinate cap does not remove it.
For \(d=1\), there are no off-diagonal pairs and the pairwise loop is skipped.
A scalar KSC result consequently provides no test of pairwise repair.

\subsection{Why a bounded coordinate map is not a stability proof}
\label{sec:bf-eot-hm-cap-stability}

The original standardized coordinate cap is, for even \(p\),
\begin{equation}
 f_{\mathrm{old}}(u)=\frac{u}{[1+(u/c)^p]^{1/p}},\qquad c>0.
 \label{eq:bf-eot-hm-original-cap}
\end{equation}
It bounds each mapped coordinate by \(c\), but compresses magnitudes throughout
the cloud.  With \(c=0.98\) and \(p=8\), moderate and large standardized
coordinates are pushed toward similar magnitudes.  After variance restoration,
that compression can leave too little variation in \(u^2\), and hence too
little kurtosis by \eqref{eq:bf-eot-hm-kurtosis-radii}.

The revised optional map is exactly the identity inside radius \(r\):
\begin{equation}
 f_{r,c}(u)=
 \begin{cases}
 u,& |u|\le r,\\
 \operatorname{sign}(u)\{r+w\tanh[(|u|-r)/w]\},& |u|>r,
 \end{cases}
 \qquad w=rc.
 \label{eq:bf-eot-hm-identity-core-cap}
\end{equation}
At the join, the outside value is \(\operatorname{sign}(u)r\), its first
derivative is one and its second derivative is zero.  It therefore joins the
identity with two continuous derivatives.  Outside, its derivative is
\(\operatorname{sech}^2[(|u|-r)/w]\), between zero and one, and
\(\lvert f_{r,c}(u)\rvert<r(1+c)\).  The September candidate uses \(r=8\)
and \(c=0.98\), giving a pre-restoration bound of 15.84.  These are recorded
candidate settings, not universal settings for another model or tuning scope.
The map is followed by centering, covariance restoration and mapping to
\((\mu_w,P_w)\); those operations are part of the reset.

There are three reasons why this bound alone cannot prevent a higher-moment
fit from going wrong.  First, the cap is applied \emph{after} the fitting
iterations, so it cannot repair an earlier invalid solve or factorization.
Second, a bounded displacement need not improve a nonlinear moment residual,
especially after rewhitening.  Third, whitening can amplify a cloud whose
smallest covariance eigenvalue has collapsed.  A bound before whitening is
not a bound after whitening.  Even well inside radius eight, a single
standardized coordinate of magnitude eight contributes \(8^4/1008\simeq4.06\)
to a 1008-particle fourth moment.  Finiteness alone is therefore a weak test
of higher-moment stability.

\subsection{A guarded finite fit with a protected fallback}
\label{sec:bf-eot-hm-guarded-fit}

The safety repair retains the scaled LM solve, both displacement caps and the
final coordinate cap.  It adds a test of each proposed step \emph{after}
rewhitening and a separate test of the complete capped reset.  The objective
must be specified because marginal and mixed moments can compete.  For a
selected target entry \(t_j\), let \(m_j(U)\) denote its empirical moment and
let \(a_j=\max(1,|t_j|)\).  Stack the marginal third and fourth moments and,
when pairwise correction is configured, the selected entries of
\eqref{eq:bf-eot-hm-pair-targets}.  The checked discrepancy is
\begin{equation}
 \mathcal L(U)=\sum_{j\in\mathcal J}
        \left(\frac{m_j(U)-t_j}{a_j}\right)^2.
 \label{eq:bf-eot-hm-guard-loss}
\end{equation}
This scale gives absolute error near a zero target and relative error for
large targets.  Symmetric co-kurtosis entries occur twice in the all-pairs
sum.  The normalization is an explicit safety objective, not a proof that
every moment has equal scientific importance.  In particular, a decrease in
\(\mathcal L\) can accompany an increase in one individual kurtosis residual;
the individual residual families must remain visible.  The marginal LM and
pairwise proposals need not be descent directions for this combined objective.
The acceptance test, rather than that assumption, decides whether to use them.

For a current standardized cloud \(U\), define its empirical covariance \(C\)
and a capped proposal displacement \(D\).  Try the fixed sequence
\(\alpha\in\{1,1/2,\ldots,1/128\}\), stopping at the first accepted trial.
For \(V_\alpha=U+\alpha D\), let \(C_\alpha\) be its centered empirical
covariance.  Before factoring it, require finite values and tangents, a valid
input covariance, and
\begin{equation}
 \tfrac14 C\preceq C_\alpha\preceq\tfrac94 C.
 \label{eq:bf-eot-hm-guard-covariance}
\end{equation}
The positive-semidefinite inequalities are checked by symmetric eigenvalues.
An inadmissible trial does not enter the Cholesky branch.  A valid trial is
centered and whitened, using the same normalization convention as its original
marginal or pairwise update, to obtain \(\mathcal W(V_\alpha)\).  Accept it only
if the resulting value, analytical tangent and loss are finite and
\begin{equation}
 \mathcal L(\mathcal W(V_\alpha))
       \le\mathcal L(U)+\tau_{\mathrm{fp}}(\mathcal L(U)),\qquad
 \tau_{\mathrm{fp}}(l)=32\epsilon_{\mathrm{mach}}(1+|l|).
 \label{eq:bf-eot-hm-guard-acceptance}
\end{equation}
If all eight trials fail, retain \(U\) and its incoming tangent.  This is a
bounded search, with no assertion of eventual acceptance or convergence.
Repeated rejection is reported, rather than concealed by a zero-loss claim.

To interpret the covariance envelope, put
\(B=C^{-1/2}C_\alpha C^{-1/2}\).  Equation
\eqref{eq:bf-eot-hm-guard-covariance} implies
\(\lambda(B)\in[1/4,9/4]\).  Consequently
\(\lVert B^{-1/2}\rVert_2\le2\) and
\(\lVert B^{1/2}\rVert_2\le3/2\).  These bounds concern relative whitening
in the current covariance metric; they do not bound the source covariance,
all accumulated derivatives, or the final physical-state coordinates.
The implementation uses a roundoff margin
\(32\epsilon_{\mathrm{mach}}\operatorname{tr}(C)/d\) in the matrix tests and
requires the input minimum eigenvalue to exceed that margin.  It also requires
the trial minimum eigenvalue to exceed
\(32\epsilon_{\mathrm{mach}}\operatorname{tr}(C_\alpha)/d\): otherwise a
Loewner roundoff allowance could admit a singular trial from an input already
near that allowance.  The displayed
exact-arithmetic bounds should therefore be read with that numerical allowance.

The last check includes operations outside the fitting loop.  Let
\(\mathcal F(U)\) mean applying the chosen coordinate cap, recentering,
restoring covariance and mapping to the weighted source scale.  Compute both
\(X_{\mathrm{fit}}=\mathcal F(U_{\mathrm{fit}})\) and
\(X_0=\mathcal F(U_0)\), where \(U_0\) is the initial standardized raw
Contract~E cloud and has undergone no moment fitting.  Evaluate their moments
in the same whitening convention used for reporting.  Return the fitted cloud
only if it is valid and
\begin{equation}
 \mathcal L(\mathcal W(X_{\mathrm{fit}}))
 \le \mathcal L(\mathcal W(X_0))+
             \tau_{\mathrm{fp}}(\mathcal L(\mathcal W(X_0))).
 \label{eq:bf-eot-hm-complete-map-check}
\end{equation}
Otherwise return \(X_0\), together with its correctly recomputed total tangent.
The baseline must itself be valid; an invalid baseline makes the route invalid.
Both alternatives use the \emph{same} coordinate cap and restoration.  Thus
rejecting a fit does not bypass the tail protection.  In this optional safety
mode, a zero-iteration call deliberately evaluates that protected baseline;
the historical zero-iteration, safety-disabled call remains a no-op.

The exact covariance proposition above assumes un-ridged whitening.  The
numerical implementation also uses, where declared, the relative ridge
\(Q_\delta=\operatorname{sym}(Q)+\delta\operatorname{tr}(Q)I/d\), with
\(\delta=10^{-12}\).  If \(RR^\top=Q_\delta\), then the empirical covariance
after whitening is
\[
 R^{-1}QR^{-\top}
 =I-\delta\operatorname{tr}(Q)R^{-1}R^{-\top}/d.
\]
It is not exactly the identity.  The ridge scale and its tangent belong to the
finite numerical program, and covariance residuals are checked numerically.
A small coefficient alone does not imply a small effect when \(Q\) is badly
conditioned.  The step guard checks the fit; it does not rescue invalid source
moments or an invalid upstream Contract~E factorization.

The implementation records the trial count, rejected-step count, smallest
accepted step, input-validity flag, largest accepted loss increase, protected
baseline loss, final loss and final-fallback flag.  It is optional through
\texttt{moment\_safety=True} in the shared marginal/pairwise routine and the
canonical filtering consumer.  Full projected-cumulant correction is a distinct
path and is explicitly rejected with this option, rather than silently omitted
from its safety objective.  The richer residual design and identity-core cap
also remain explicit settings; none has been promoted to a universal default.


\paragraph{Bounded numerical evidence.}
The October 2026 safety comparison used fresh fixed fixtures in dimensions
1, 2, 3 and 10, with healthy, skewed, mixture, outlier, concentrated-weight,
nearly collinear and repeated-axis cases.  All 112 GPU/XLA cases passed the
finite-value and declared loss checks; a 14-case CPU reference also passed.
The guard rejected at least one complete proposed step in 23 GPU cases and
used a reduced step without a complete rejection in another 33 cases.
These counts describe the evaluated fixtures, not probabilities of failure.
The guard was less accurate in the declared combined moment loss than at least
one simple comparator in 26 of the 112 GPU cases, including mixture,
concentrated-weight and repeated-axis regimes.  This vetoes an accuracy-based
promotion.  It does not contradict the narrower safety comparison with the
identically protected no-fit cloud.

A subsequent edge-case audit added the explicit positive trial-eigenvalue
check described above.  Its focused CPU regression passed, and eight further
GPU/XLA checks of the final implementation covered active tail capping,
healthy-update agreement and analytical derivatives.  The largest FP64
finite-difference tangent discrepancy in those checks was
\(1.18\times10^{-10}\).  Healthy GPU particle values were unchanged; their
largest tangent difference was \(4.34\times10^{-19}\), so bitwise equality of
all GPU tangents is not asserted.  The broad comparison's largest relative
covariance residual was \(1.10\times10^{-3}\) in FP32/TF32 and
\(1.32\times10^{-15}\) in FP64.  These measured distinctions are why the
exact covariance proposition must not be presented as an unconditional
floating-point identity.  The complete comparisons, source versions, failures,
commands and limitations are recorded in
\path{docs/benchmarks/ledh-moment-safety-results-20261001.md}.
The original historical explosive input was not recovered and replayed;
fresh stress fixtures cannot establish that every such failure is resolved.

\subsection{Complete algorithm and score statement}

For the LEDH--PFPF--OT filter, the following sequence specifies one complete
time step.  It is essential to distinguish the proposal stage from the reset:
the bootstrap weights in \eqref{eq:bf-eot-hm-likelihood}--
\eqref{eq:bf-eot-hm-weights} illustrate the reset's input, whereas an LEDH
proposal requires the importance correction in Chapter~\ref{ch:bf-dpf-pfpf}.
Using observation likelihoods alone after a general particle-flow proposal
would compute different weights.
\begin{enumerate}[label=(\arabic*)]
 \item Carry each particle, its UKF posterior covariance, its weight and their
       analytical parameter tangents from the preceding step.  Predict its
       local Gaussian, generate its fixed innovation, and execute the declared
       finite LEDH flow using the particle's own covariance lifecycle.
       Apply the proposal-density and flow-Jacobian importance correction;
       accumulate the log normalizer and normalize the resulting weights.
       Update the local UKF covariance and its total tangent.  These local
       covariances are distinct from the global empirical reset covariance.
 \item Regard the corrected weighted cloud as
       \(\{x_{t,i}^{-},\alpha_{t,i}\}_{i=1}^N\).  Compute its mean and
       covariance, diagonal shape targets, and any selected pairwise targets.
       Keep the source and weight derivatives in all these quantities.
 \item Run the fixed finite Sinkhorn and balancing schedule, form the
       row-normalized barycenters in \eqref{eq:bf-eot-hm-barycentric}, and check
       the declared marginal and validity diagnostics.  The streaming
       implementation uses the repository's exact-divisor chunk policy.
 \item Inject the frozen residual design through
       \eqref{eq:bf-eot-hm-residual} and apply the realized-covariance
       Contract~E restoration \eqref{eq:bf-eot-hm-contract-e}.  The revised
       option uses \eqref{eq:bf-eot-hm-normal-quantile-design}; it does not
       change the OT weights or claim to reproduce the full source law.
       Standardize this raw reset to obtain \(U_0\).
 \item For each of the declared \(S\) marginal iterations, compute the moment
       directions, scaled LM coefficients and capped displacement in
       \eqref{eq:bf-eot-hm-directions}--\eqref{eq:bf-eot-hm-smooth-trust-cap}.
       With safety enabled, apply the covariance and post-whitening tests
       \eqref{eq:bf-eot-hm-guard-covariance}--
       \eqref{eq:bf-eot-hm-guard-acceptance}; reduce or reject the displacement
       as specified.  With safety disabled, execute the original capped
       update and whitening.
 \item For each of the declared \(P\) pairwise iterations, form
       \eqref{eq:bf-eot-hm-pair-residuals}, project and cap its direction by
       \eqref{eq:bf-eot-hm-pair-projection}--\eqref{eq:bf-eot-hm-pair-cap},
       and apply the same optional guard and rewhitening.  Skip this loop
       when \(d=1\) or \(P=0\).
 \item Apply the declared coordinate map, recenter, restore covariance and
       map back to the source mean and scale.  With safety enabled, compare
       this complete result with the identically capped and restored no-fit
       cloud using \eqref{eq:bf-eot-hm-complete-map-check}.  Select both value
       and total tangent from the accepted alternative and retain the
       diagnostics.  Fail closed if the protected baseline is invalid.
 \item Carry the selected particles with equal weights, together with the
       reset route's per-particle covariance carry and its analytical
       derivative, to the next observation.  Continue the score recursion
       through the same branch as the value recursion.
\end{enumerate}
All iteration counts, masks, caps, damping controls, residual permutations and
halving schedules are fixed within an evaluation.  Scope-specific offline
tuning, when needed for an accuracy claim, precedes the untouched claim run;
this safety mechanism is not permission to retune controls inside HMC.
The particle count satisfies the chosen residual-design rule: the implemented
replicated-axis and normal-quantile designs require \(2d\mid N\), while an
exactly replicated positive GenUT design requires representable masses.
The correction retains all \(N\) rows.  Marginal and pairwise value calculations
use \(O(Nd+d^2)\) additional storage beyond the OT workspace, with additional
storage for the requested tangent directions.  Pairwise products and covariance
checks add \(O(Nd^2+d^3)\) work per trial; eight trials bound the work, not the
moment error.

The analytical tangent must follow all dependencies of this finite algorithm.
For example, writing a dot for a parameter-direction derivative,
\begin{align}
 \dot s_a&=\sum_i\dot\alpha_i z_{ia}^3
          +3\sum_i\alpha_i z_{ia}^2\dot z_{ia},\\
 \dot K^\star_{ab}&=\sum_i\dot\alpha_i z_{ia}^2z_{ib}^2
       +2\sum_i\alpha_i(z_{ia}\dot z_{ia}z_{ib}^2
                        +z_{ia}^2z_{ib}\dot z_{ib}).
 \label{eq:bf-eot-hm-target-total-tangents}
\end{align}
The weight terms cannot be dropped.  If \(CC^\top=P\), put
\(E=C^{-1}\dot P C^{-\top}\) and let \(\mathcal T(E)\) retain its strict
lower triangle and half its diagonal.  Then
\begin{equation}
 \dot C=C\mathcal T(E),\qquad
 \dot u=(\dot x-\dot\mu)C^{-\top}
           -u\dot C^\top C^{-\top}.
 \label{eq:bf-eot-hm-whitening-total-tangent}
\end{equation}
These formulas also differentiate a declared relative ridge through its trace
scale.  Differentiate the LM linear system by solving
\(M\dot y=\dot b-\dot M y\), where \(b=A^\top d\); differentiate the
pairwise matrix products and smooth caps by the ordinary chain rule.
The identity-core cap multiplies a tangent by one in its core and by the
\(\operatorname{sech}^2\) factor in its tail.

On a locally unchanged acceptance branch, the selected \(\alpha\) is a fixed
member of the halving schedule, so
\(\dot V_\alpha=\dot U+\alpha\dot D\).  A rejected step carries \(\dot U\),
and a final fallback carries the total derivative of \(X_0\), rather than a
zero tangent.  The eigenvalue comparisons and acceptance decisions select a
branch; they are not differentiated as if they were smooth functions.
Consequently branch boundaries may have a discontinuous derivative or value.
A finite-difference check must verify that its two perturbations select the
same branches, and local agreement is not a global smoothness or HMC theorem.

\begin{proposition}[Finite-route total-score correctness]
\label{prop:bf-eot-hm-score}
Fix the innovation rows and ordering, residual design, finite Sinkhorn and
balance counts, \(\epsilon\), OT ridge, \(\rho\), LM damping \(\lambda\),
column-scale floor \(\varepsilon_s\), trust radius \(\tau\), marginal and
pairwise counts, masks, coordinate cap, and safety schedule.
Assume every required covariance is strictly positive definite, every row mass
is positive, the maximum branch in the cost scale and all safety-acceptance
and final-fallback branches are fixed locally, and every validity condition holds.  If the recursive tangent includes every operation
listed above, then its accumulated score is the total derivative with respect
to \(\theta\) of the same finite scalar
\(\widehat L^N(\theta)=\sum_t\widehat\ell_t^N(\theta)\) executed by the value
program.
\end{proposition}

\begin{proof}
On the stated branch, the transition and likelihood adapters, normalization,
finite Sinkhorn recurrences, row quotient, moment maps, matrix products,
Cholesky factors, triangular solves, column scales, \(2\times2\) damped solves,
smooth trust and coordinate caps, selected pairwise operations, and finite
rewhitening maps are differentiable on this branch.  The ordinary chain rule therefore
applies to their finite composition.  The implemented tangent includes the
dependence of the normalized weights, OT kernel and scalings, target moments,
directions \(q_3,q_4\), Jacobians \(J_a\), column scales \(H_a\), scaled
systems \(M_a\), damped coefficients \(c_a\), row cap factors \(g_i\), and both
standardization maps.  Induction over time then shows that the carried
particle tangent is the derivative of the carried particle value and that the
sum of increment tangents is \(d\widehat L^N/d\theta\).

This proves correctness for the executed finite approximation.  It does not
prove that \(\widehat L^N\) equals the exact filtering log likelihood or that
its derivative equals the exact filtering score.  Those are model-specific
statistical approximation questions tested against available references.
\end{proof}

\paragraph{Correctness boundary.}
The propositions establish the OT quotient semantics, the Contract~E ridged
identity, exact post-correction empirical mean and covariance under the un-ridged
full-rank assumptions, and the local total derivative of one fixed finite
program.  Declared ridges have the covariance discrepancy derived above;
safety selection gives its stated loss bound, not global differentiability.  They do not establish global
higher-moment feasibility, exact mixed moments, unbiased likelihood, exact
posterior score, superiority over Contract~E, HMC readiness, or a default
change.  The model campaign below can reject this candidate or keep it viable;
it cannot by itself prove broad nonlinear-model validity.


\subsection{What the actual likelihoods and scores show}
\label{sec:bf-eot-hm-observed-results}

The downstream question is whether preserving the weighted cloud's shape
brings the accumulated likelihood and score closer to the filtering reference.
For the scalar KSC example, the full seven-mixture likelihood and its two
parameter derivatives can be evaluated by a separately refined quadrature.
The September comparison used this reference at
\(\theta=(\gamma,\log\beta)=(1.5,0)\), two untouched datasets of length 120,
and eight particle designs per method.  It ran in FP64 with GPU/XLA and TF32
disabled.  Tables~\ref{tab:bf-eot-hm-ksc-120-a} and
\ref{tab:bf-eot-hm-ksc-120-b} report the actual quantities, including the
Gaussian Kalman approximation as a simple comparator.  The revised arm uses
the richer residual, identity-core cap and controls selected for that scope;
the comparison does not isolate the causal effect of each change.
The full-mixture reference is an oracle for the declared seven-mixture target,
rather than for a different stochastic-volatility observation law.
The numerical values, refinement checks and component uncertainties are
preserved in the comparison record.\footnote{\path{docs/plans/artifacts/sqmc-ksc-reset-repair-20260929/attempt-01/analysis-01/comparisons.md}}

\begin{table}[htbp]
\centering
\small
\setlength{\tabcolsep}{5pt}
\begin{tabular}{lrrr}
\toprule
Method & Log likelihood & \(\gamma\) score & \(\log\beta\) score\\
\midrule
Seven-mixture reference & -270.702896202 & -0.508659556 & -1.291005193\\
Gaussian Kalman heuristic & -275.542306325 & -0.813750759 & -0.892859070\\
IID / original & -271.316778239 & -0.432267129 & -1.085171066\\
IID / revised & -270.879543409 & -0.633906248 & -1.217855750\\
Inverse CDF / original & -271.239752570 & -0.375158888 & -1.089143021\\
Inverse CDF / revised & -270.821584886 & -0.580225487 & -1.269873809\\
Permutation / original & -271.238543975 & -0.371483139 & -1.087564459\\
Permutation / revised & -270.823316362 & -0.577390069 & -1.269882173\\
Ablation / original & -271.239262624 & -0.368794285 & -1.087245064\\
Ablation / revised & -270.823316362 & -0.577390069 & -1.269882173\\
\bottomrule
\end{tabular}
\caption{KSC at \(T=120\), dataset A (seed 243001).
Particle entries are means of eight independent designs with \(N=1008\).
The revised arm uses the September design/cap repair and its selected controls;
it does not include the October safety guard.}
\label{tab:bf-eot-hm-ksc-120-a}
\end{table}

\begin{table}[htbp]
\centering
\small
\setlength{\tabcolsep}{5pt}
\begin{tabular}{lrrr}
\toprule
Method & Log likelihood & \(\gamma\) score & \(\log\beta\) score\\
\midrule
Seven-mixture reference & -284.105359942 & 2.358935308 & 0.299061235\\
Gaussian Kalman heuristic & -289.119307428 & 2.448686517 & 0.655680321\\
IID / original & -284.231222413 & 2.856512226 & 0.299130287\\
IID / revised & -284.459443680 & 2.338869846 & 0.347294048\\
Inverse CDF / original & -284.160036239 & 3.035272629 & 0.331255873\\
Inverse CDF / revised & -284.225769251 & 2.368748070 & 0.350155526\\
Permutation / original & -284.159942453 & 3.026332424 & 0.330404077\\
Permutation / revised & -284.227196844 & 2.366669388 & 0.350044821\\
Ablation / original & -284.158477299 & 3.046262376 & 0.333190588\\
Ablation / revised & -284.227196844 & 2.366669388 & 0.350044821\\
\bottomrule
\end{tabular}
\caption{KSC at \(T=120\), dataset B (seed 243002).
Particle entries are means of eight independent designs with \(N=1008\).
The revised arm uses the September design/cap repair and its selected controls;
it does not include the October safety guard.}
\label{tab:bf-eot-hm-ksc-120-b}
\end{table}

The inverse-CDF route illustrates the distinction between likelihood and score
accuracy.  On dataset A its mean absolute log-likelihood error across designs
fell from \(0.536856\) to \(0.163187\); on dataset B it rose from
\(0.125318\) to \(0.213643\).  Its mean score-vector error norm fell from
\(0.329639\) to \(0.144456\) on A and from \(0.681811\) to \(0.146093\) on B.
These are means of per-design errors, which differ from the error of a mean
reported in the tables.  The paired 95\% bootstrap intervals for revised
minus original score-error norm were \([-0.291356,-0.084593]\) and
\([-0.654304,-0.399304]\), respectively.  They support a conditional reduction
on these two datasets; they are exploratory intervals across several examined
cells, without a multiplicity adjustment.

The uncertainty in each score coordinate also matters.  For the revised
inverse-CDF arm, the standard errors of the mean \(\gamma\) and
\(\log\beta\) scores were \((0.053617,0.012488)\) on A and
\((0.055220,0.026190)\) on B.  Thus a visible difference in the last digits is
not itself evidence of a persistent bias.  Conversely, the near equality of
the original IID \(\log\beta\) score and reference on B does not establish
that every coordinate improved after repair.  The aggregate score norm and
the individual coordinates answer different questions.  These results make
the revised mechanism worth testing further, while its failed strict accuracy
screen and the dataset-B likelihood deterioration prevent a general accuracy
or default-readiness conclusion.

To see why a likelihood offset and a score discrepancy need not move together,
write the error of a fixed-design finite filter as
\begin{equation}
 e_N(\theta)=\widehat L^N(\theta)-L(\theta).
 \label{eq:bf-eot-hm-likelihood-error}
\end{equation}
Here \(L\) is the exact log likelihood for the declared model.  At a parameter
where both functions are differentiable and the numerical selection branches
are locally fixed, subtracting their derivatives gives
\begin{equation}
 \widehat s_N(\theta)-s(\theta)=\nabla_\theta e_N(\theta).
 \label{eq:bf-eot-hm-score-error}
\end{equation}
A constant likelihood offset can therefore coexist with an accurate score;
a small likelihood error at one parameter can coexist with a steep error
gradient.  Likewise the error in a log-likelihood ratio between two parameters
is \(e_N(\theta')-e_N(\theta)\), not either offset separately.  Assessing a
parameter neighborhood, both score coordinates and branch changes is
necessary before using these observations to argue for gradient-based
inference.  The tables at one parameter do not establish HMC readiness.

There are also two distinct fourth-moment questions.  Fitting can bring a reset
closer to the weighted particle cloud's fourth moment.  That target may itself
differ from the filtering distribution's fourth moment.  Let \(m_4^{R}\),
\(m_4^{W}\) and \(m_4^\star\) denote the reset, weighted-cloud and reference
fourth moments, all expressed in one common centering and scaling convention.
Then
\[
 m_4^{R}-m_4^\star
   =(m_4^{R}-m_4^{W})+(m_4^{W}-m_4^\star).
\]
The proposed correction addresses the first term.  The second requires an
independent filtering reference or a better source cloud.  Comparing likelihoods
and scores with their oracle directly tests the downstream approximation, but
does not by itself identify which fourth-moment term caused an error.

\subsection{The first nonlinear comparisons}

Predator--prey and spatial SIR test different weaknesses of the repair.
Predator--prey has two state coordinates and nonlinear coupled dynamics, so
pairwise moments already have a substantive role.  Its six score coordinates
are the physical parameters \(r,K,a,s,u,v\) used by the declared transition.
The spatial SIR case has 18 state coordinates, nine observed infectious
coordinates, and three parameters: the log infection-rate scale, log
removal-rate scale and log observation-noise scale.  Its larger collection of
mixed moments tests whether fitting dependence remains numerically useful as
dimension grows.

A master tester now sends both models through the same finite LEDH executor,
per-particle UKF covariance recursion, Contract~E reset and analytical total
score.  It compares covariance-only, original capped, richer marginal-only,
richer marginal-and-pairwise, and guarded marginal-and-pairwise configurations.
The last pair isolates the optional guard; the marginal-only comparison
isolates the additional mixed-moment correction.  Likelihoods, every score
coordinate, replication uncertainty, invalid evaluations and reset diagnostics
are retained separately for each model, dataset, parameter point, particle
count and ancestry rule.\footnote{\path{docs/benchmarks/ledh-nonlinear-master.md}}

Fresh synthetic observations follow the declared initial Gaussian law and a
transition before each observation.  This timing matters: a bootstrap filter
that observes the initial state before its first transition estimates a
different likelihood for the same numerical observations.  The independent
reference used here starts from the same initial law and saved observations
as LEDH, then transitions before every observation.  It also retains the same
parameter values, including their original floating-point representation.

The reference is a Gaussian bootstrap particle filter with an analytical
Fisher-identity score.  Write \(x_{0:T}\) for a complete latent trajectory,
\(f_\theta\) for its transition density and \(g_\theta\) for the observation
density.  Because the initial law in these examples is parameter independent,
regular differentiation under the integral gives
\[
 \nabla_\theta\log p_\theta(y_{1:T})
 =\mathbb E_\theta\!\left[
   \sum_{t=1}^T\{\nabla_\theta\log f_\theta(x_t\mid x_{t-1})
              +\nabla_\theta\log g_\theta(y_t\mid x_t)\}
       \,\middle|\,y_{1:T}\right].
\]
Each particle carries this accumulated complete-data score through resampling;
the final weighted average estimates the conditional expectation.  State
values are held fixed when differentiating the complete-data densities, and
their Gaussian covariance derivatives are included.  This is an approximation
to the observed-data score, whereas LEDH differentiates its own finite
numerical likelihood.  Their finite-particle scores need not coincide.

The reference passed fixed-state finite-difference checks for both nonlinear
models and a separate linear-Gaussian comparison with an exact Kalman
likelihood and score.  Four independent replications were then evaluated at
8,192, 32,768, 131,072 and 524,288 particles.  The reported reference likelihood
is the logarithm of the mean likelihood estimate; its score averages the
Fisher estimates with likelihood weights.  Delete-one-replication jackknife
standard errors describe Monte Carlo variation.  They do not certify that
finite-particle bias has vanished, so this reference is not called an exact
nonlinear oracle.

The first comparison used \(T=20\) and \(N=1008\) for LEDH, with inherited,
untuned numerical controls.  The two predator--prey datasets used seeds
260401 and 260402; the SIR dataset used 260401.  LEDH ran on the GPU with
XLA, FP32 and TF32.  Predator--prey returned finite likelihoods and all six
score coordinates in 70 evaluations across the five correction choices and
the tested ancestry rules.  Sixty-nine also passed the separate trace-score
check.  Table~\ref{tab:bf-nonlinear-first-likelihoods} gives the IID means,
which make the difference between the two models visible.

\begin{table}[htbp]
\centering
\small
\begin{tabular}{lrrr}
\toprule
Model / dataset & Covariance only & Guarded pairwise & Reference (MCSE)\\
\midrule
Predator--prey, 260401 & $-97.43870$ & $-97.35035$ & $-97.34297\;(0.00657)$\\
Predator--prey, 260402 & $-99.29759$ & $-99.28713$ & $-99.31152\;(0.00366)$\\
SIR, 260401, FP64 & $-975.57354$ & $-1041.35394$ & $-678.07747\;(0.01486)$\\
\bottomrule
\end{tabular}
\caption{Actual log likelihoods at \(T=20\).  LEDH means use two IID designs
on dataset 260401 and four on predator--prey dataset 260402.  The SIR row is a
same-data FP64 precision diagnostic; all ten initial FP32 SIR evaluations
failed.  The SIR guarded mean includes one finite evaluation whose separate
trace-score check failed.  References use four replications at 524,288
particles.  These small-sample means do not establish a ranking.}
\label{tab:bf-nonlinear-first-likelihoods}
\end{table}

For predator--prey dataset 260401, the guarded score mean in the order
\((r,K,a,s,u,v)\) was
\[
 (-33.52772,\;-0.54696,\;0.01964,\;4.56826,\;-8.54557,\;10.63757),
\]
while the reference was
\[
 (-32.90524,\;-0.54394,\;0.02053,\;4.48734,\;-8.73240,\;10.86545).
\]
The reference coordinate MCSEs were
\[
 (0.07550,0.00339,0.0000942,0.00827,0.02200,0.02770).
\]
For dataset 260402, the corresponding guarded and reference score means were
\begin{align*}
 &(-21.21510,\;0.03513,\;0.01890,\;1.52016,\;-6.60116,\;8.16064),\\
 &(-21.14020,\;0.03931,\;0.01868,\;1.50221,\;-6.53110,\;8.07440).
\end{align*}
The reference coordinate MCSEs there were
\[
 (0.09711,0.00155,0.00000666,0.00506,0.00271,0.00270).
\]
The full per-design scores, LEDH replication standard errors and conditional
comparisons with covariance-only, original and marginal-only correction are
retained in the result record.\footnote{\path{docs/benchmarks/ledh-nonlinear-execution-results-20261002.md}}
The values support further predator--prey validation, but two or four design
replications at a single parameter point do not establish that the repair
improves the observed-data score.

SIR exposes two separate problems.  In the localized FP32 run, the first
failure occurred at observation 4 in the Contract~E reset.  Ancestry, UKF
prediction, UKF update and model-callback validity checks still passed; flow
states and derivatives and likelihood weights were finite.  Covariance-only
reset states remained finite but their derivatives did not.  The guarded
configuration lost finite reset states and derivatives, so the higher-moment
guard no longer had a valid input to protect.  These observations locate the
failure in the reset but do not yet identify its exact factorization or
prove a remedy.

Replaying the exact FP32 observations and random inputs in FP64 gave finite
SIR likelihoods and scores.  Accuracy still failed conspicuously relative to
the independent reference.  The two guarded likelihoods were \(-1041.82982\)
and \(-1040.87806\), with score vectors
\begin{align*}
 &(-165.24294,\;-186.13243,\;963.39536),\\
 &(3230.85495,\;-1581.97484,\;936.41325).
\end{align*}
In the log-scale parameter order given above, the reference score was
\[
 (106.31321,-65.96031,5.58635),
\]
with coordinate MCSEs \((5.53617,2.42038,0.03538)\).
Thus restoring finite arithmetic did not restore filtering accuracy.  The
hundreds of log-likelihood units separating SIR's LEDH and reference estimates
are far larger than the measured reference replication noise.

Two finite evaluations also retained trace-score consistency vetoes: original
FP32 predator--prey differed by \(6.47\times10^{-4}\) in its first score
coordinate, and one guarded FP64 SIR evaluation differed by
\(1.71\times10^{-7}\).  The recorded tolerances were not relaxed.  Together
with the absence of scope-specific calibration and untouched validation,
these findings prevent admission or a default-readiness conclusion.
The next SIR step is to diagnose the failing reset factorization and its
analytical derivative, then evaluate any protection for non-harm before
calibrating the approximation.  It is not evidence that fitting fourth
moments should be abandoned: the failure occurs in a common reset stage,
and the independent likelihood and score comparison remains necessary even
when that stage is numerically healthy.

\section{A Zhao--Cui squared-TT moment teacher for Contract E}
\label{sec:bf-eot-zhao-cui-moment-teacher}

The empirical targets in
\eqref{eq:bf-eot-hm-target-moments} preserve higher moments already present in
the weighted particle cloud.  They cannot correct Monte Carlo error in those
same empirical moments.  This section defines a different, opt-in
construction: run a deterministic squared tensor-train (TT) density recursion
in parallel with the particle filter, use the TT density only to compute
standardized shape moments, and give those moments to the higher-moment
Contract~E correction.  The construction avoids a dense tensor grid.

The source-grounded operations are the adjacent-density recursion, squared-TT
nonnegative representation, and paired-core marginal contractions in
Algorithms~1--2 and Proposition~2 of \citet{zhao2024ttsequential}.  The author
implementation constructs the adjacent approximation in
\code{third\_party/audit/zhao\_cui\_tensor\_ssm\_p10/source/models/full\_sol.m},
lines 21--136, and performs squared-TT marginal recursions in
\code{third\_party/audit/zhao\_cui\_tensor\_ssm\_p10/source/deep-tensor.dev/src/@TTSIRT/marginalise.m},
lines 25--85.  Using the resulting moments as Contract~E targets is not an
operation claimed by Zhao and Cui.  The assembled route is therefore
\code{extension\_or\_invention}; it must not be called a source-faithful
Zhao--Cui filter.

\subsection{Parallel TT filtering recursion}

The teacher is evaluated conditionally at the current HMC parameter
\(\theta\).  The parameter is not an additional random coordinate and is never
integrated out.  Only latent/static variables explicitly belonging to the
filtering state enter the TT chart.  In the state-only case the adjacent order
is \((x_t,x_{t-1})\).  If \(\vartheta\) denotes a genuine carried latent
variable, its coordinate order and marginalization must be declared by the
model.  Given the previous conditional approximation
\(\widehat\pi_{t-1}^{\rm TT}(x_{t-1},\vartheta;\theta)\), define
\begin{equation}
 q_t(x_t,\vartheta,x_{t-1})
 =\widehat\pi_{t-1}^{\rm TT}(x_{t-1},\vartheta)
   p_\theta(x_t\mid x_{t-1},\vartheta)
   p_\theta(y_t\mid x_t,\vartheta).
 \label{eq:bf-eot-tt-teacher-adjacent-target}
\end{equation}
Here \(\vartheta\) must not be identified with the HMC parameter \(\theta\).
Every factor in \(q_t\) is evaluated at the fixed \(\theta\) supplied by the
outer likelihood program.
On a declared product-measure coordinate chart
\(r=(r_1,\ldots,r_m)\), first pull the target back to the reference measure.
If \(x=\Psi_t(r)\) is the complete chart, define
\begin{equation}
 \bar q_t(r)=\frac{q_t(\Psi_t(r))\left|\det D\Psi_t(r)\right|}
 {\omega_\nu(r)},\qquad
 \omega_\nu(r)=\frac{d\nu}{dr}.
 \label{eq:bf-eot-tt-teacher-chart-pullback}
\end{equation}
The Jacobian factor and, when \(\nu\) is not Lebesgue measure, the inverse
reference density belong to the Radon--Nikodym target and its tangent; omitting
either computes moments under the wrong measure.  Define
\begin{equation}
 c_t=\max_{\text{fit rows}}\log\bar q_t(r),\qquad
 s_t(r)=\exp\!\left\{\tfrac12[\log\bar q_t(r)-c_t]\right\}.
 \label{eq:bf-eot-tt-teacher-square-root-target}
\end{equation}
The fixed-branch fitter minimizes its declared weighted ridge objective
against \(s_t\), not against \(\bar q_t\).  If \(h_t\) denotes the fitted
square-root TT, form
\begin{equation}
 \lambda_t\geq0,\qquad
 \tau_t=e^{-c_t}\lambda_t,\qquad
 \widetilde q_t(r)=h_t(r)^2+\tau_tq_{0,t}(r),\qquad
 Z_t^{\rm TT}=\int\widetilde q_t(r)\,d\nu(r),
 \quad \nu=\bigotimes_{k=1}^m\nu_k.
 \label{eq:bf-eot-tt-teacher-squared-density}
\end{equation}
Here \(q_{0,t}\geq0\), \(\lambda_t\) is the defensive weight in the
unscaled pulled-back target, and \(\tau_t\) is its coefficient in scaled TT
units.  Thus an exact fit gives
\(\widetilde q_t=e^{-c_t}(\bar q_t+\lambda_tq_{0,t})\).  If instead a route
declares \(\tau_t\) directly and keeps it fixed, it defines a different
scale-shift-dependent defensive mixture and must carry that semantics in its
route identity.  The normalized teacher density is
\(\widehat q_t^{\rm TT}=\widetilde q_t/Z_t^{\rm TT}\).  Marginalizing the
coordinates for \(\vartheta\) and \(x_{t-1}\) gives
\begin{equation}
 \widehat\pi_t^{\rm TT}(x_t)
 =\int\widehat q_t^{\rm TT}(x_t,\vartheta,x_{t-1})
       \,d\vartheta\,dx_{t-1}.
 \label{eq:bf-eot-tt-teacher-filter-marginal}
\end{equation}
The recursion also carries the \((x_t,\vartheta)\) marginal needed at the next
time.  It is therefore a parallel approximate density filter, not a one-time
fit to current particle moments.  Fitting it to empirical particle moments
 would destroy the intended independence of the teacher.  Fitting \(h_t\) to
 \(\bar q_t\), rather than to
 \eqref{eq:bf-eot-tt-teacher-square-root-target}, would approximately represent
 \(\bar q_t^2\) and is wrong relative to this density target.

All ranks, bases, coordinate order, fitting designs, sweep counts, ridge and
floor controls, and discrete solver branches are fixed for one score-bearing
program.  ``Fixed branch'' means these choices do not change with \(\theta\);
 it does not mean that the fitted cores are numerically constant in \(\theta\).

\begin{proposition}[Square-root fit target and conditional parameter semantics]
\label{prop:bf-eot-tt-teacher-square-root-target}
Assume \(\bar q_t(r)>0\) on the declared fit rows and let \(c_t\) be independent
of the fit-row index.  If \(h_t(r)=s_t(r)\) exactly, then
\(h_t(r)^2=\exp(-c_t)\bar q_t(r)\).  Thus, with \(\tau_t=0\), normalization
of \(h_t^2\) gives the normalized conditional target.  Treating \(\theta\) as
a chart coordinate and integrating it instead changes the probability measure
and defines a different algorithm.
\end{proposition}

\begin{proof}
Squaring \eqref{eq:bf-eot-tt-teacher-square-root-target} gives
\(s_t(r)^2=\exp(-c_t)\bar q_t(r)\).  The positive scalar
\(\exp(-c_t)\) cancels from the normalizing quotient.  The conditional
statement follows because \(\theta\) is fixed in every factor of \(q_t\);
integrating a variable requires specifying a measure for it and therefore
defines a different joint target.  For the unscaled defensive weight in
\eqref{eq:bf-eot-tt-teacher-squared-density}, both components have the common
factor \(e^{-c_t}\), which cancels.  A coefficient \(\tau_t\) held fixed while
\(c_t\) changes does not share this factor and therefore changes the
normalized density.
\end{proof}

\subsection{Exact fixed-ALS fit and its forward derivative}

Let the frozen fit rows be \(r^{(n)}\), with positive weights \(w_n\), and
write \(s_n=s_t(r^{(n)})\).  A realized sweep schedule is a fixed sequence of
active axes \(k_1,\ldots,k_U\).  Immediately before update \(u\), contract all
non-active cores at the fit rows to obtain the design
\(A_u\in\mathbb R^{M\times p_u}\).  Vectorizing the active core as \(c_u\),
the declared ridge update is
\begin{align}
 H_u&=A_u^\top W A_u+\rho I,
 &b_u&=A_u^\top W s,
 &H_uc_u&=b_u,
 \label{eq:bf-eot-tt-teacher-als-update}
\end{align}
where \(W=\operatorname{diag}(w_1,\ldots,w_M)\).  The implementation may use
objective-preserving column scaling and an augmented least-squares solve, but
the returned unscaled coefficient vector must solve
\eqref{eq:bf-eot-tt-teacher-als-update}.  Any condition-number or floor branch
is part of the frozen finite route.

For one parameter direction, let dots denote total directional derivatives of
the state immediately before update \(u\).  Because \(A_u\) depends on every
non-active core, its derivative must be built from their current tangents.  The
exact update JVP is
\begin{align}
 \dot H_u
 &=\dot A_u^\top W A_u+A_u^\top W\dot A_u
   +A_u^\top\dot W A_u+\dot\rho I,\\
 \dot b_u
 &=\dot A_u^\top Ws+A_u^\top\dot W s+A_u^\top W\dot s,\\
 H_u\dot c_u
 &=\dot b_u-\dot H_uc_u.
 \label{eq:bf-eot-tt-teacher-als-jvp}
\end{align}
After each value update, both \(c_u\) and \(\dot c_u\) replace the active
core and its tangent before constructing update \(u+1\).  Recomputing every
\(\dot c_u\) from only the terminal fitted cores, or setting
\(\dot A_u=0\) after another core has acquired a tangent, is a partial
derivative and is wrong relative to the fixed-ALS program.

On a locally fixed maximizing row \(n_t^*\), the scale and target tangents are
\begin{equation}
 \dot c_t=\dot{\log\bar q}_t(r^{(n_t^*)}),
 \qquad
 \dot s_n=\tfrac12s_n
 [\dot{\log\bar q}_t(r^{(n)})-\dot c_t].
 \label{eq:bf-eot-tt-teacher-square-root-jvp}
\end{equation}
The derivative of \(\log\bar q_t\) includes the analytical model score, any
moving-chart/Jacobian term, and, for \(t>0\), the tangent of the normalized
previous TT marginal.  Fixed reference rows and weights have
\(\dot W=0\); a parameter-dependent design or ridge must instead include the
corresponding terms above.  The current reference replay fixes the reference
rows, basis, weights, and ridge, hence uses \(\dot W=0\) and \(\dot\rho=0\),
while retaining \(\dot A_u\) through all non-active core tangents.  It cannot
be reused for a moving-basis, moving-row, or parameter-dependent-ridge route
without adding those missing derivatives.
For the scale-consistent defensive convention in
\eqref{eq:bf-eot-tt-teacher-squared-density},
\begin{equation}
 \dot\tau_t=e^{-c_t}(\dot\lambda_t-\lambda_t\dot c_t)
 \label{eq:bf-eot-tt-teacher-defensive-scale-jvp}
\end{equation}
including the zero-weight case.  Omitting the
\(-\tau_t\dot c_t\) contribution changes the derivative of the represented
defensive mixture.

The scale shift is a numerical conditioning device, not a likelihood
increment in the hybrid particle route.  In particular, its tangent must enter
the square-root fit target through
\eqref{eq:bf-eot-tt-teacher-square-root-jvp}, but neither \(c_t\) nor
\(\log Z_t^{\rm TT}\) is added to the particle likelihood.  The latter is the
normalizer of an auxiliary approximate density, whereas the particle
increment is defined by the normalized particle weights before the reset.

\begin{proposition}[Sequential fixed-ALS JVP]
\label{prop:bf-eot-tt-teacher-als-jvp}
Assume the update schedule and all discrete branches are locally fixed and
every \(H_u\) is nonsingular.  Starting from the derivative of the declared
initial cores, sequential application of
\eqref{eq:bf-eot-tt-teacher-als-jvp} returns the total directional derivative
of every terminal fitted core produced by
\eqref{eq:bf-eot-tt-teacher-als-update}.
\end{proposition}

\begin{proof}
At the first update, differentiating \(H_1c_1=b_1\) gives
\(H_1\dot c_1=\dot b_1-\dot H_1c_1\), hence the stated formula.  Suppose all
core values and tangents before update \(u\) equal those of the finite value
program.  Product-rule contraction of the non-active cores gives the total
\(\dot A_u\); differentiating \(H_uc_u=b_u\) gives the stated
\(\dot c_u\) by a linear solve with the same \(H_u\).  Replacing the active
core and tangent establishes the induction
hypothesis for update \(u+1\).  Induction over the finite schedule proves the
claim.
\end{proof}

\subsection{Exact observable contractions for a fixed squared TT}

Write the square-root TT in a basis \(\phi_{k\ell}\) as
\begin{equation}
 h(r)=G_1(r_1)\cdots G_m(r_m),\qquad
 G_k(r_k)=\sum_{\ell=1}^{n_k}C_{k\ell}\phi_{k\ell}(r_k),
 \label{eq:bf-eot-tt-teacher-functional-tt}
\end{equation}
where \(C_{k\ell}\in\mathbb R^{R_{k-1}\times R_k}\) and
\(R_0=R_m=1\).  For a separable observable
\(g(r)=\prod_kg_k(r_k)\), define
\begin{equation}
 M_k[g_k]_{\ell s}
 =\int\phi_{k\ell}(u)\phi_{ks}(u)g_k(u)\,d\nu_k(u)
 \label{eq:bf-eot-tt-teacher-operator-matrix}
\end{equation}
and the paired-core transfer
\begin{equation}
 A_k[g_k]_{(a,A),(b,B)}
 =\sum_{\ell,s}C_k[a,\ell,b]C_k[A,s,B]M_k[g_k]_{\ell s}.
 \label{eq:bf-eot-tt-teacher-paired-transfer}
\end{equation}

\begin{proposition}[Squared-TT observable contraction]
\label{prop:bf-eot-tt-teacher-observable-contraction}
For \eqref{eq:bf-eot-tt-teacher-functional-tt} and any integrable separable
\(g\),
\begin{equation}
 \int g(r)h(r)^2\,d\nu(r)
 =e_1^\top A_1[g_1]\cdots A_m[g_m]e_1.
 \label{eq:bf-eot-tt-teacher-observable-contraction}
\end{equation}
Consequently
\begin{equation}
 \mathbb E_{\widehat q^{\rm TT}}[g]
 =\frac{e_1^\top A_1[g_1]\cdots A_m[g_m]e_1
       +\tau\int gq_0\,d\nu}
      {e_1^\top A_1[1]\cdots A_m[1]e_1
       +\tau\int q_0\,d\nu}.
 \label{eq:bf-eot-tt-teacher-normalized-observable}
\end{equation}
\end{proposition}

\begin{proof}
Expand both copies of \(h\) in their basis and TT-rank indices.  Product
measure and separability allow Fubini's theorem to split every basis-pair
integral by coordinate.  The integral at coordinate \(k\) is exactly
\(M_k[g_k]\); summing the two TT rank paths gives
\eqref{eq:bf-eot-tt-teacher-paired-transfer}.  Left-to-right contraction proves
\eqref{eq:bf-eot-tt-teacher-observable-contraction}.  Linearity adds the
defensive term, and division by the same expression with \(g=1\) proves
\eqref{eq:bf-eot-tt-teacher-normalized-observable}.
\end{proof}

It is wrong to contract unsquared TT cores with integral vectors: moments of
\(h^2\) require paired cores and
\eqref{eq:bf-eot-tt-teacher-operator-matrix}.  It is also wrong to omit the
\(\tau q_0\) numerator term when \(\tau>0\).

For coordinate monomials choose \(g_k(u)=u^{\alpha_k}\).  For an affine
current-state chart
\begin{equation}
 x_t=a_t+B_tr_t,
 \label{eq:bf-eot-tt-teacher-affine-chart}
\end{equation}
the contracted reference mean \(m_r\) and covariance \(P_r\) give
\begin{equation}
 \mu_t^{\rm TT}=a_t+B_tm_r,\qquad
 P_t^{\rm TT}=B_tP_rB_t^\top.
 \label{eq:bf-eot-tt-teacher-affine-moments}
\end{equation}

\begin{proposition}[Mean and covariance of the TT teacher]
\label{prop:bf-eot-tt-teacher-affine-moments}
If \(Z_t^{\rm TT}>0\), the monomial contractions and
\eqref{eq:bf-eot-tt-teacher-affine-moments} equal the mean and covariance of
the normalized finite density \(\widehat q_t^{\rm TT}\) under the declared
affine chart.
\end{proposition}

\begin{proof}
Apply Proposition~\ref{prop:bf-eot-tt-teacher-observable-contraction} with
\(g=r_i\) and \(g=r_ir_j\).  These give \(m_r\) and
\(P_r=\mathbb E[rr^\top]-m_rm_r^\top\).  Standard affine moment identities
then give \eqref{eq:bf-eot-tt-teacher-affine-moments}.
\end{proof}
\subsection{Low-rank contractions for whitened higher moments}

Let \(L_t^{\rm TT}(L_t^{\rm TT})^\top=P_t^{\rm TT}\) and define
\begin{equation}
 z=(L_t^{\rm TT})^{-1}(x_t-\mu_t^{\rm TT}).
 \label{eq:bf-eot-tt-teacher-whitening}
\end{equation}
Each component is an affine form in the reference coordinates,
\(z_i=c_i+\sum_ka_{ik}r_k\).  Directly forming every raw third- and
fourth-order tensor would require \(O(d^3)\) and \(O(d^4)\) outputs.  That is
not the proposed high-dimensional algorithm.

For nonnegative integers \(p,q\), maintain a bivariate polynomial state
\(v_{ab}=S^aT^b\), \(0\leq a\leq p\), \(0\leq b\leq q\).  When coordinate
\(r_k\) adds \(\alpha_kr_k\) to \(S\) and \(\beta_kr_k\) to \(T\), the exact
local update is
\begin{equation}
 (S+\alpha_kr_k)^a(T+\beta_kr_k)^b
 =\sum_{u=0}^a\sum_{v=0}^b
   \binom{a}{u}\binom{b}{v}S^uT^v
   \alpha_k^{a-u}\beta_k^{b-v}r_k^{a-u+b-v}.
 \label{eq:bf-eot-tt-teacher-polynomial-automaton}
\end{equation}
Initialize with \(S=c_i,T=c_j\), apply
\eqref{eq:bf-eot-tt-teacher-polynomial-automaton} over the coordinates, and
select state \((p,q)\).  Replacing each \(r_k^s\) by the operator matrix
\(M_k[r_k^s]\) couples this fixed-size automaton to the paired-core
contraction.  Its auxiliary dimension is at most \((p+1)(q+1)\), hence at most
9 for the pairwise moments below.

\begin{proposition}[Affine-form mixed moments without a dense moment tensor]
\label{prop:bf-eot-tt-teacher-affine-form-moments}
Assume \(M_k[r_k^s]\) is exact for \(0\leq s\leq p+q\).  The coupled
paired-core/polynomial-automaton contraction returns
\begin{equation}
 \int (c_i+a_i^\top r)^p(c_j+a_j^\top r)^q h(r)^2\,d\nu(r)
 \label{eq:bf-eot-tt-teacher-affine-form-integral}
\end{equation}
exactly for the represented TT \(h\), without constructing a dense third- or
fourth-order moment tensor.
\end{proposition}

\begin{proof}
Equation~\eqref{eq:bf-eot-tt-teacher-polynomial-automaton} is the binomial
theorem in the two accumulated affine forms.  Induction over \(k\) shows that
the selected final state is exactly the product in
\eqref{eq:bf-eot-tt-teacher-affine-form-integral}.  Expanding that automaton
path and the two copies of the TT, then applying Fubini coordinate by
coordinate, replaces each local monomial by its exact operator matrix.  The
remaining rank contractions are precisely the stated integral.
\end{proof}

The teacher supplies
\begin{equation}
 \gamma_i^{\rm TT}=\mathbb E[z_i^3],\quad
 \kappa_i^{\rm TT}=\mathbb E[z_i^4],\quad
 C_{ij}^{(3),\rm TT}=\mathbb E[z_i^2z_j],\quad
 C_{ij}^{(4),\rm TT}=\mathbb E[z_i^2z_j^2].
 \label{eq:bf-eot-tt-teacher-shape-targets}
\end{equation}
The two pairwise quantities have different index semantics.  The co-skew term
\(C_{ij}^{(3),\rm TT}=\mathbb E[z_i^2z_j]\) is ordered: requesting \((i,j)\)
does not request \((j,i)\), which would instead be
\(\mathbb E[z_j^2z_i]\).  The co-kurtosis term is symmetric, so one requested
unordered pair \(\{i,j\}\) supplies both matrix entries
\(C_{ij}^{(4),\rm TT}=C_{ji}^{(4),\rm TT}\).  The implementation therefore
maintains separate structural masks for ordered co-skew and symmetric
co-kurtosis entries; uncomputed entries are never silently interpreted as zero
targets.  Only declared pairs \((i,j)\in\mathcal E\) need be contracted.  All pairs
give \(O(d^2)\) outputs and at least \(O(d^2)\) work even though no dense
fourth-order tensor is formed.  For a sparse interaction graph or other
predeclared structural set with \(|\mathcal E|=O(d)\), the number of requested
shape moments remains linear in \(d\).  Pair selection is fixed outside HMC
and cannot be tuned on the claim path.

\subsection{Composition with Contract E}

The default composition is deliberately hybrid.  Let \(\mu_t^w,P_t^w\) be
the likelihood-weighted particle moments from
\eqref{eq:bf-eot-hm-target-moments}.  Contract~E and the final affine map keep
\(\mu_t^w,P_t^w\) as their first-two-moment targets.  Only the standardized
shape targets in \eqref{eq:bf-eot-tt-teacher-shape-targets} replace their
empirical counterparts in the bounded correction
\eqref{eq:bf-eot-hm-coefficients}--\eqref{eq:bf-eot-hm-map-back}.  This keeps
the current Contract~E mean/covariance claim intact while giving the shape
correction information not computed from the same particle cloud.

Using \(\mu_t^{\rm TT},P_t^{\rm TT}\) themselves as reset targets is a
different algorithm.  It may be tested as an ablation, but it does not
preserve the particle weighted covariance and must not be silently substituted
for the hybrid route.

\begin{proposition}[Hybrid first-two-moment preservation]
\label{prop:bf-eot-tt-teacher-hybrid-preservation}
Assume \(P_t^w\) and every correction covariance are positive definite.  If
the bounded correction is rewhitened after each iteration and mapped back with
the Cholesky factor of \(P_t^w\), the final equally weighted cloud has mean
\(\mu_t^w\) and covariance \(P_t^w\), irrespective of the TT shape targets.
\end{proposition}

\begin{proof}
The TT targets affect only the finite correction directions and coefficients.
After the final correction, Proposition~\ref{prop:bf-eot-hm-rewhitening}
gives zero mean and identity covariance in standardized coordinates.  Mapping
back by \(\mu_t^w\) and the Cholesky factor of \(P_t^w\) proves the result.
\end{proof}

The higher moments are not guaranteed to match
\eqref{eq:bf-eot-tt-teacher-shape-targets}: the correction uses damped finite
steps, moment constraints can conflict, and rewhitening changes higher
moments.  Post-correction residuals are mandatory diagnostics.

\subsection{Likelihood, score, and derivative statements}

The particle likelihood increment \(\widehat\ell_t^N\) is computed before the
TT-driven reset.  Adding the teacher after that increment leaves the already
computed scalar at time \(t\) unchanged.  The reset changes the cloud carried
to \(t+1\), so it can and generally will change future approximate increments.
There is no theorem that it preserves the full filtering distribution, the
complete finite-horizon particle likelihood, or the exact model likelihood.

For a core direction \(\dot C_k\), operator direction \(\dot M_k\), and
transfer \(A_k=A_k(C_k,M_k)\), the product rule gives
\begin{equation}
 \dot A_k
 =A_k(\dot C_k,C_k,M_k)+A_k(C_k,\dot C_k,M_k)
  +A_k(C_k,C_k,\dot M_k),
 \label{eq:bf-eot-tt-teacher-transfer-jvp}
\end{equation}
and, for \(H_g=e_1^\top\prod_kA_k e_1\),
\begin{equation}
 \dot H_g
 =\sum_{j=1}^m e_1^\top
   A_1\cdots A_{j-1}\dot A_jA_{j+1}\cdots A_m e_1.
 \label{eq:bf-eot-tt-teacher-contraction-jvp}
\end{equation}
If \(N_g=H_g+\tau Q_{0,g}\), \(Z=N_1\), then
\begin{equation}
 \frac{d}{d\theta}\mathbb E_{\widehat q^{\rm TT}}[g]
 =\frac{\dot N_gZ-N_g\dot Z}{Z^2},\qquad
 \dot N_g=\dot H_g+\dot\tau Q_{0,g}+\tau\dot Q_{0,g}.
 \label{eq:bf-eot-tt-teacher-normalized-jvp}
\end{equation}
The same product rule applies to polynomial-automaton transfers, including
chart Jacobian, Cholesky, and affine-form coefficient derivatives.  For a
non-affine chart, a whitened state component is not generally a polynomial in
the reference coordinates.  The polynomial automaton is then valid only when
the transformed observable is represented in a declared finite basis;
otherwise fixed quadrature must construct the operator matrices, and its error
is an additional approximation that must be measured.

\begin{proposition}[Total score of the fixed TT-teacher finite program]
\label{prop:bf-eot-tt-teacher-total-score}
Fix the TT bases, ranks, fitting designs, sweep counts, regularization,
coordinate ordering, pair set \(\mathcal E\), particle innovations, OT
controls, and all discrete branches.  Assume every finite fit solve and
Cholesky factor is nonsingular near \(\theta\).  If the recursive tangent
includes (i) every fitted-core tangent, including dependence through the
previous TT marginal; (ii) every observable, normalizer, chart, whitening, and
defensive-density term in
\eqref{eq:bf-eot-tt-teacher-transfer-jvp}--\eqref{eq:bf-eot-tt-teacher-normalized-jvp};
and (iii) the complete particle, OT, Contract~E, and correction tangent, then
the accumulated score equals the total derivative of the same executed finite
likelihood program.
\end{proposition}

\begin{proof}
On the fixed branch, each TT fit is a finite composition of smooth target
evaluations and nonsingular fixed-design solves.  Equations
\eqref{eq:bf-eot-tt-teacher-transfer-jvp}--\eqref{eq:bf-eot-tt-teacher-normalized-jvp}
are the product and quotient rules for the exact finite contractions.
Cholesky and triangular-solve JVPs differentiate whitening, and
Proposition~\ref{prop:bf-eot-hm-score} differentiates the particle/OT/reset
composition.  Induction over time includes the dependence of the next TT fit
on the previous TT marginal and of the next particle increment on the reset
cloud.  The chain rule gives the derivative of the same finite scalar returned
by the value program.
\end{proof}

A tangent through only the final TT contraction, treating fitted cores as
constants, is a partial derivative and is wrong relative to this total-score
claim.  A pre-fitted teacher frozen across \(\theta\) defines a different
finite approximation; a zero target tangent can be correct for that explicitly
frozen approximation but is not the score of the recursively refitted teacher.
Runtime rank selection, stochastic fitting choices, finite-difference
production scores, and an autodiff-only score route are ineligible for the
intended HMC path.

\subsection{Complete algorithm and complexity boundary}

Before the first observation, freeze the basis family, ranks, fit and holdout
rows, fit weights, coordinate charts, ALS sweep schedule, ridge and
conditioning policy, defensive density, ordered co-skew mask, symmetric
co-kurtosis mask, particle innovations, OT controls, Contract~E controls, and
bounded-correction controls.  These are part of the finite-program identity;
none may be reselected as \(\theta\) changes.  Initialize the TT teacher from
the declared initial-state density and initialize the particle cloud by the
ordinary fixed particle route.

For each observation time \(t\), execute the following ordered program:
\begin{enumerate}[label=(\arabic*)]
 \item Hold the outer HMC parameter \(\theta\) fixed.  Evaluate the conditional
       adjacent log target from
       \eqref{eq:bf-eot-tt-teacher-adjacent-target}.  At \(t=0\), replace the
       previous TT marginal by the declared initial density.  At \(t>0\),
       evaluate the normalized carried TT marginal and its quotient-rule
       tangent.  Never add \(\theta\) as a TT integration coordinate.
 \item Pull the target back through the complete chart using
       \eqref{eq:bf-eot-tt-teacher-chart-pullback}.  Include the chart Jacobian,
       inverse reference density, and all their tangents.  Select the fixed
       local maximizing row and construct \(s_t,\dot s_t\) using
       \eqref{eq:bf-eot-tt-teacher-square-root-target} and
       \eqref{eq:bf-eot-tt-teacher-square-root-jvp}.
 \item Starting from the declared value and tangent cores, replay every ALS
       update in its frozen order.  At update \(u\), rebuild \(A_u\) and
       \(\dot A_u\) from the current cores, solve
       \eqref{eq:bf-eot-tt-teacher-als-update}, solve
       \eqref{eq:bf-eot-tt-teacher-als-jvp}, then immediately replace both the
       active core and its tangent.  Veto a non-finite value, branch change,
       excessive condition number, or failed primal/JVP algebraic residual.
 \item Form \(\widetilde q_t=h_t^2+\tau_tq_{0,t}\), its normalizer, and their
       tangents.  Contract the previous-state coordinates and any genuine
       carried latent coordinates, retaining the normalized
       \((x_t,\vartheta)\) TT marginal and its tangent for time \(t+1\).
 \item By paired-core contractions, compute the teacher mean and covariance.
       Differentiate the covariance, Cholesky factor, and triangular solves to
       obtain the whitening coefficients and their tangents.  Contract
       diagonal skew/kurtosis and only the declared ordered co-skew and
       symmetric co-kurtosis entries, including the defensive numerator,
       normalizer, and quotient tangents.
 \item Independently propagate the particle rows with the frozen innovations,
       evaluate their observation weights, and record
       \(\widehat\ell_t^N\), the normalized weights, and their analytical
       tangents.  This particle scalar, not \(c_t+\log Z_t^{\rm TT}\), is the
       hybrid likelihood increment.
 \item Execute the declared OT quotient and canonical Contract~E--Chol reset.
       This produces an equally weighted cloud with the particle weighted mean
       and covariance; propagate the total particle/weight/transport tangent.
 \item Apply the fixed number of bounded higher-moment correction steps using
       the TT shape values and tangents.  Rewhiten after every step, map back
       with the particle weighted mean and covariance, and record every
       post-correction target residual.  The residuals are diagnostics; the
       method does not claim exact higher-moment matching.
 \item Carry the corrected equal-weight particle cloud and tangent, together
       with the normalized TT \((x_t,\vartheta)\) marginal and tangent, to
       time \(t+1\).  After the final observation return
       \(\sum_t\widehat\ell_t^N\) and the sum of its analytically propagated
       tangents.
\end{enumerate}

This order makes the two probability calculations explicit.  The particle
lane owns the finite likelihood value and score.  The TT lane owns only an
auxiliary normalized filtering approximation used to define shape targets.
The two lanes meet at the reset after the current particle increment has been
recorded; therefore the TT normalizer cannot be substituted for the particle
increment without changing the algorithm.

For maximum basis size \(n\), maximum TT rank \(R\), and observable-automaton
dimension \(\rho\leq9\), the density stores \(O(mnR^2)\) coefficients and a
contraction has working state \(O(\rho R^2)\).  A conservative dense local
contraction costs up to \(O(m\rho^2n^2R^4)\), with lower cost available from
contraction ordering and orthogonal bases.  Mean/covariance output costs
\(O(d^2)\).  Selected moments cost at least \(O(|\mathcal E|)\) outputs.  These
costs do not depend on particle count \(N\), but the per-step TT fit can
dominate them and TT rank can grow sharply with nonlinear coupling.  No
dimension-independent rank bound is claimed for the models of interest.

\paragraph{Correctness and promotion boundary.}
The propositions prove exact moments and their JVPs for a fixed represented
squared-TT density, and exact first-two particle moments after the hybrid
reset.  They do not prove that the TT density approximates the true filter,
that its rank stays bounded, that the correction reaches its TT targets, that
likelihood or score bias decreases, or that one TT fit per step is affordable
inside a high-dimensional nonlinear HMC evaluation.  Dense-quadrature parity,
fixed-branch derivative parity, LGSSM/Kalman checks, nonlinear value/score
no-regression, and GPU/XLA memory/runtime evidence remain separate gates.

The current variable-rank TT containers and the reference fixed-ALS replay in
\code{bayesfilter/highdim/zhao\_cui\_moment\_teacher\_als.py} store cores as a
Python tuple, so a
straight implementation of the products above uses a setup-static Python loop
that TensorFlow traces into an unrolled graph.  That is sufficient for an
independent TensorFlow mechanics reference, but it is ineligible for the
BayesFilter XLA runtime under the no-Python-loop policy.  Candidate runtime
promotion additionally requires a padded/masked or otherwise batch-native core
representation whose contraction loop is a TensorFlow control-flow operation.
The reference modules implement the conditional square-root target, sequential
fixed-ALS value/JVP equations, normalized carried marginal, and TT shape-target
JVP in a setup-static recursion.  A separate setup-static composition accepts
the particle likelihood increment and tangent, runs the existing OT and
Contract~E--Chol JVP, and applies the TT shape targets while returning the
particle scalar unchanged.  These references are not a model-scale filter and
do not yet implement the graph-native form.

For the padded candidate, let every core be embedded in
\(\widetilde G_k\in\mathbb R^{R\times B\times R}\), and let a fixed mask
\(P_k\) select the entries of the variable-rank core.  All inactive entries of
\(\widetilde G_k\) and \(\dot{\widetilde G}_k\) are zero.  At an update of
axis \(k\), the padded design \(\widetilde A_k\) has zero columns outside
\(P_k\).  More precisely, for \(M\) prepared rows and
\(p=R^2B\), let
\(\widetilde A_k\in\mathbb R^{M\times p}\),
\(W=\operatorname{diag}(w_1,\ldots,w_M)\in\mathbb R^{M\times M}\),
\(s\in\mathbb R^M\), and let
\(P_k\in\mathbb R^{p\times p}\) be the diagonal zero--one projector
onto the active core entries.  Then
\(\widetilde A_k=\widetilde A_kP_k\).  For a fixed \(\rho>0\), the
padded ridge equations split as
\begin{equation}
 \left(\widetilde A_k^\top W\widetilde A_k+\rho I\right)
 \widetilde g_k=\widetilde A_k^\top Ws,
 \qquad
 \widetilde g_k=P_k\widetilde g_k\in\mathbb R^p,
 \label{eq:bf-eot-tt-teacher-padded-als}
\end{equation}
because every inactive equation is \(\rho\widetilde g_{k,j}=0\), while the
active block is exactly equation~\eqref{eq:bf-eot-tt-teacher-als-update}.
Thus padding does not alter the represented update when the ridge is strictly
positive.  The zero-ridge case is not part of this candidate: choosing a
pseudoinverse convention for the inactive block would define a different
finite numerical program.

Let \(S_k\) be the fixed-branch positive diagonal column scale computed from
the primal padded design.  The candidate solves
\begin{equation}
 \left[S_k^{-1}\left(\widetilde A_k^\top W\widetilde A_k+\rho I\right)
 S_k^{-1}\right]z_k
 =S_k^{-1}\widetilde A_k^\top Ws,
 \qquad \widetilde g_k=S_k^{-1}z_k .
 \label{eq:bf-eot-tt-teacher-padded-scaled-solve}
\end{equation}
This is an invertible change of coordinates, not a changed objective.  The JVP
is still equation~\eqref{eq:bf-eot-tt-teacher-als-jvp} embedded in the padded
space; differentiating the auxiliary scale is unnecessary because the claimed
coefficient is characterized by the unscaled normal equations.  The graph
propagates the left and right environments and their tangents with TensorFlow
control flow and consumes a fixed update schedule.  It separately gates the
condition of the scaled primal system and the unscaled derivative system, and
records both algebraic residuals.  Masks, ranks, bases, rows, schedule, ridge,
and fit count are fixed before tracing.  Hence this construction removes the
Python core/sweep loops without introducing parameter-dependent branch or rank
selection; it does not by itself establish that the fitted TT is an accurate
filtering approximation.

\FloatBarrier
